%% ========================================================================
%% Jiuzhang Suanshu (九章算術) -- Volumes 7-9
%% BILINGUAL English-Chinese Edition
%% With Commentary by Liu Hui (劉徽)
%% ========================================================================
\documentclass[11pt,a4paper]{article}

\usepackage{fontspec}
\usepackage{polyglossia}
\usepackage{amsmath,amssymb,amsthm}
\usepackage{geometry}
\usepackage{graphicx}
\usepackage{xcolor}
\usepackage{fancyhdr}
\usepackage{array,booktabs,longtable}
\usepackage{hyperref}
\usepackage{xeCJK}
\usepackage{verbatim}
\usepackage{fancyvrb}
\usepackage{titlesec}
\usepackage{enumitem}
\usepackage{multicol}
\usepackage{paracol}
\usepackage{tcolorbox}
\tcbuselibrary{skins,breakable}

\geometry{
  paperwidth=297mm,paperheight=210mm,
  left=15mm,right=15mm,top=22mm,bottom=22mm,
  headheight=14pt,footskip=30pt
}

\setmainlanguage{english}
\setmainfont{TeXGyreTermes}
\setsansfont{TeXGyreHeros}
\setmonofont[Scale=MatchLowercase]{TeXGyreCursor}

\setCJKmainfont[AutoFakeBold=true,AutoFakeSlant=true]{Noto Sans CJK SC}
\setCJKsansfont[AutoFakeBold=true,AutoFakeSlant=true]{Noto Sans CJK SC}
\setCJKmonofont[AutoFakeBold=true,AutoFakeSlant=true]{Noto Sans CJK SC}

\definecolor{titlecolor}{RGB}{80,10,10}
\definecolor{sectioncolor}{RGB}{120,20,20}
\definecolor{volcolor}{RGB}{45,55,72}
\definecolor{commentarycolor}{RGB}{20,80,20}
\definecolor{methodbg}{RGB}{240,240,255}
\definecolor{commentbg}{RGB}{240,255,240}
\definecolor{quotecolor}{RGB}{60,40,20}
\definecolor{cnlabelfg}{RGB}{120,20,20}
\definecolor{enlabelfg}{RGB}{20,80,120}

\columnsep=0.4cm

\newcommand{\cnwen}{\textcolor{cnlabelfg}{\textbf{問：}}}
\newcommand{\cnda}{\textbf{答曰：}}
\newcommand{\cnshu}{\textcolor{cnlabelfg}{\textbf{術曰：}}}
\newcommand{\cnzhu}{\textcolor{commentarycolor}{\textit{\textbf{劉徽注：}}}}
\newcommand{\enwen}{\textcolor{enlabelfg}{\textbf{Q:}}}
\newcommand{\enda}{\textcolor{enlabelfg}{\textbf{A:}}}
\newcommand{\enshu}{\textcolor{enlabelfg}{\textbf{Method:}}}
\newcommand{\enzhu}{\textcolor{enlabelfg}{\textit{\textbf{Liu Hui's Commentary:}}}}

\newenvironment{mathdisplay}{\switchcolumn[0]*}{\switchcolumn[0]}

\pagestyle{fancy}
\fancyhf{}
\fancyhead[L]{\small\textsc{Jiuzhang Suanshu Volumes 7--9 -- Bilingual Edition}}
\fancyhead[R]{\small\thepage}
\renewcommand{\headrulewidth}{0.4pt}
\renewcommand{\footrulewidth}{0pt}

\hypersetup{
  colorlinks=true, linkcolor=sectioncolor, citecolor=sectioncolor, urlcolor=sectioncolor,
  pdfencoding=auto,
  pdftitle={Jiuzhang Suanshu Volumes 7-9 Bilingual Edition},
  pdfauthor={Anonymous; Commentary by Liu Hui}
}

\setcounter{tocdepth}{2}
\setcounter{secnumdepth}{2}

\newcommand{\vs}{\vspace{0.5em}}
\begin{document}

\thispagestyle{empty}
\begin{center}
\vspace*{1.5cm}
{\fontsize{36}{44}\selectfont\color{titlecolor} 九章算術}
\vspace{0.5cm}
{\fontsize{18}{22}\selectfont\color{volcolor} （卷七至卷九）}
\vspace{0.3cm}
{\LARGE\itshape Jiuzhang Suanshu --- Volumes 7--9}
\vspace{0.4cm}
{\large\itshape Nine Chapters on the Mathematical Art}
\vspace{2cm}
{\LARGE\bfseries\color{sectioncolor} BILINGUAL EDITION}
\vspace{0.3cm}
{\large\itshape 中英文對照版}
\vspace{2cm}
\begin{tabular}{ll}
\textbf{Volume 7:} & 盈不足 Yingbuzu --- Excess and Deficit \\
\textbf{Volume 8:} & 方程 Fangcheng --- Rectangular Arrays \\
\textbf{Volume 9:} & 勾股 Gougu --- Right Triangles
\end{tabular}
\vspace{2cm}
{\Large 魏\quad 劉徽\quad 注}
\vspace{0.4cm}
{\large Commentary by Liu Hui (c.~263 CE)}
\vspace{1.5cm}
\rule{0.4\textwidth}{1pt}
\vspace{0.8cm}
{\small\textsc{Left: Original Chinese \quad | \quad Right: English Translation}}
\vspace{0.3cm}
{\small Typeset in LaTeX\ with XeLaTeX and xeCJK}
\end{center}

\newpage
\tableofcontents
\newpage

%% ========================================================================
%% PREFACE
%% ========================================================================
\section*{序\quad /\quad Preface}
\addcontentsline{toc}{section}{序 / Preface}

\begin{paracol}{2}
\switchcolumn[0]
\begin{quote}\color{quotecolor}\small
\textit{The \textbf{Jiuzhang Suanshu} (九章算術) stands as one of the supreme achievements of ancient mathematics. Compiled over centuries and finalized around the Han dynasty, it presents 246 problems across nine chapters. Liu Hui's commentary (c.~263 CE) transformed this practical handbook into a work of deep mathematical insight.}

\textit{The three volumes presented here represent the mathematical crown jewels: the algorithmic brilliance of Yingbuzu (double false position), the algebraic sophistication of Fangcheng (Gaussian elimination), and the geometric elegance of Gougu (the Pythagorean theorem).}
\end{quote}

\switchcolumn[1]
\begin{quote}\color{quotecolor}\small
\textit{\textbf{《九章算術》}是古代數學最偉大的成就之一。經數百年編纂，約成書於漢代，全書九章共二百四十六題。劉徽注約成於西元263年，將這部實用手冊提升為具有深刻數學洞察力的經典。}

\textit{本書所呈現的三卷，乃《九章》之數學瑰寶：盈不足之算法精妙、方程之代數造詣、勾股之幾何優雅。}
\end{quote}
\end{paracol}

\vspace{0.5cm}

\begin{paracol}{2}
\switchcolumn[0]
\cnshu 劉徽序曰：昔在包犧氏始畫八卦，以通神明之德，以類萬物之情，作九九之術，以合六爻之變。暨於黃帝神農，其為功也，各有稱焉。周公制禮而有九數，九數之流，則九章是矣。

\cnzhu 徽幼習九章，長再詳覽。觀陰陽之割裂，總算術之根源，探賾之暇，遂悟其意。是以敢竭頑魯，采其所見，為之作注。事物有對，故立正負之率；錯綜無形，故開方程之術。庶幾令文約而義博，事省而功多。

\switchcolumn[1]
\enshu Liu Hui's Preface: In ancient times, Paoxi (Fuxi) first drew the Eight Trigrams to communicate the virtue of the divine spirits and to classify the nature of all things. He created the art of multiplication tables to match the transformations of the six lines. By the time of the Yellow Emperor and Shennong, each had their own accomplishments to be celebrated. The Duke of Zhou established the rites and created the Nine Numbers; from the Nine Numbers flows the Nine Chapters.

\enzhu In my youth I studied the Nine Chapters, and as I grew older I reviewed them again in detail. Observing the division of yin and yang, and summarizing the root sources of mathematical methods, in my leisure moments of deep inquiry I gradually comprehended their meaning. Therefore I dared, with all my humble capacity, to collect what I have observed and write this commentary. Things have their opposites, so I establish the rates of positive and negative; when complexity has no visible form, I open the method of rectangular arrays. My hope is to make the text concise yet the meaning broad, the effort small yet the achievement great.
\end{paracol}

\newpage

%% ========================================================================
%% VOLUME 7: YINGBUZU (盈不足)
%% ========================================================================
\newpage
\section*{卷第七\quad 盈不足\quad /\quad Volume 7\quad Excess and Deficit}
\addcontentsline{toc}{section}{卷第七 盈不足 / Volume 7 Excess and Deficit}

\begin{paracol}{2}
\switchcolumn[0]
\begin{quote}\color{quotecolor}\small
\textit{盈不足術乃《九章》中最著名之算法之一。其法以兩次假設及其盈肭推求真實之數，即「雙假設法」。此法經伊斯蘭世界傳入中世紀歐洲，被稱為 elchataym。本章共二十四題，涵蓋商業交易、幾何應用及著名之「五家共井」不定方程。}
\end{quote}
\switchcolumn[1]
\begin{quote}\color{quotecolor}\small
\textit{The method of Yingbuzu (excess and deficit) is one of the most celebrated algorithms in the Nine Chapters. It provides a technique for solving problems involving two guesses and their resulting excesses or deficits --- essentially the ``rule of double false position.'' This method was transmitted to the West via the Islamic world and medieval Europe, where it became known as the method of elchataym.}
\end{quote}
\end{paracol}

\subsection*{Method Summary / 術概}

\begin{paracol}{2}
\switchcolumn[0]
\cnshu 盈不足術曰：置所出率，盈、不足各居其下。令維乘所出率，并以為實。并盈、不足為法。實如法而一。

\cnzhu 此術以盈不足之率，求所出之實。維乘者，交錯相乘也。其術之意：假設兩端，一為盈，一為不足，以兩設交錯相乘，取其和為實，取其差之和為法，實如法而一，得所求之數。

\switchcolumn[1]
\enshu The Method of Excess and Deficit says: Place the assumed rates (guesses), with the excess and deficit placed below each respectively. Cross-multiply the assumed rates by the excess/deficit, and sum to form the dividend. Sum the excess and deficit to form the divisor. Divide the dividend by the divisor to obtain the result.

\enzhu This method uses the rates of excess and deficit to find the true value. ``Cross-multiply'' means to multiply in alternate fashion. The idea: make two assumptions, one yielding excess and one yielding deficit; cross-multiply, take their sum as dividend and the sum of differences as divisor; dividing yields the sought number.
\end{paracol}

\begin{mathdisplay}
\begin{tcolorbox}[colback=methodbg,colframe=sectioncolor,title=General Formula,fonttitle=\bfseries,boxrule=0.5pt]
For two guesses $a_1, a_2$ with excess/deficit $e_1, e_2$:
\[\text{Result} = \frac{a_1 e_2 + a_2 e_1}{e_1 + e_2}\]
\end{tcolorbox}
\end{mathdisplay}

\bigskip\hrule\bigskip

\subsection{第一題 / Problem 1: 共買物 / Buying Goods Together}

\begin{paracol}{2}
\switchcolumn[0]
\cnwen 今有共買物，人出八，盈三；人出七，不足四。問：人數、物價各幾何？

\cnda 人七，物價五十三。

\cnshu 術曰：置人出八、人出七，盈三、不足四。令維乘之：八乘四得三十二，七乘三得二十一。并之得五十三，為物價。并盈不足得七，為人數。

\cnzhu 按盈不足之術，兩設皆乖於實。假令每人出八，則多三；每人出七，則少四。維乘者，以盈乘不足之設，以不足乘盈之設，所以通而同之也。并之為實，并盈不足為法。實如法而一，得物價。

\switchcolumn[1]
\enwen Now there is a group buying goods. If each person contributes 8, there is an excess of 3; if each contributes 7, there is a deficit of 4. Question: how many people are there, and what is the price of the goods?

\enda There are 7 people; the goods cost 53.

\enshu Method: Place 8 and 7 as the two contributions, with excess 3 and deficit 4 below them. Cross-multiply: 8 times 4 gives 32; 7 times 3 gives 21. Summing gives 53, which is the price of the goods. Summing the excess and deficit gives 7, which is the number of people.

\enzhu According to the method of excess and deficit, both assumptions deviate from the truth. If each person contributes 8, there is a surplus of 3; if each contributes 7, there is a shortfall of 4. Cross-multiplication means multiplying the excess by the assumption with deficit, and multiplying the deficit by the assumption with excess, thereby harmonizing and equalizing them.
\end{paracol}

\bigskip\hrule\bigskip

\subsection{第二題 / Problem 2: 共買雞 / Buying Chickens Together}

\begin{paracol}{2}
\switchcolumn[0]
\cnwen 今有共買雞，人出九，盈十一；人出六，不足十六。問：人數、雞價各幾何？

\cnda 人九，雞價七十。

\cnshu 術曰：置人出九、人出六，盈十一、不足十六。令維乘之：九乘十六得一百四十四，六乘十一得六十六。并之得二百一十，為實。并盈不足得二十七，為法。

\cnzhu 此術與上術意同。九乘不足十六者，以多出之率乘不足之數也；六乘盈十一者，以少出之率乘盈之數也。所以維乘者，兩設不同，不可偏用，故交錯相乘以齊同之。

\switchcolumn[1]
\enwen Now there is a group buying chickens. If each person contributes 9, there is an excess of 11; if each contributes 6, there is a deficit of 16. Question: how many people are there, and what is the price of the chickens?

\enda There are 9 people; the chickens cost 70.

\enshu Method: Place 9 and 6 as contributions, with excess 11 and deficit 16. Cross-multiply: 9 times 16 gives 144; 6 times 11 gives 66. Summing gives 210 as the dividend. Summing excess and deficit gives 27 as the divisor.

\enzhu This method shares the same idea as the previous one. Nine times the deficit 16 means multiplying the larger contribution rate by the deficit amount; six times the excess 11 means multiplying the smaller contribution rate by the excess amount. The purpose of cross-multiplication is that the two assumptions differ and cannot be used one-sidedly; therefore we multiply alternately to equalize and harmonize them.
\end{paracol}

\bigskip\hrule\bigskip

\subsection{第三題 / Problem 3: 共買璡 / Buying Jade Together}

\begin{paracol}{2}
\switchcolumn[0]
\cnwen 今有共買璡，人出半，盈四；人出少半，不足三。問：人數、璡價各幾何？

\cnda 人四十二，璡價十七。

\cnshu 術曰：置人出半、人出少半，盈四、不足三。半即二分，少半即三分。令維乘之，如法求之。

\cnzhu 半者，二分之一也；少半者，三分之一也。出半則多四，出少半則少三。以分數為率，故先通分之，然後維乘求之。此術以分數入盈不足，所以見術之無所不通也。

\switchcolumn[1]
\enwen Now there is a group buying jade. If each person contributes half, there is an excess of 4; if each contributes ``less-than-half'' (one-third), there is a deficit of 3. Question: how many people are there, and what is the price of the jade?

\enda There are 42 people; the jade costs 17.

\enshu Method: Place ``half'' and ``less-than-half'' as contributions, with excess 4 and deficit 3. Half is one-half, less-than-half is one-third. Cross-multiply and solve by the method.

\enzhu ``Half'' means one-half; ``less-than-half'' means one-third. Contributing half gives excess 4; contributing one-third gives deficit 3. Since fractions are used as rates, first bring them to common denominators, then cross-multiply and solve. This method applies fractions to excess and deficit, showing that the method has no limitation in its applicability.
\end{paracol}

\bigskip\hrule\bigskip

\subsection{第四題 / Problem 4: 共買牛 / Buying Oxen Together}

\begin{paracol}{2}
\switchcolumn[0]
\cnwen 今有共買牛，七家共出一百九十，不足三百三十；九家共出二百七十，盈三十。問：家數、牛價各幾何？

\cnda 一百二十六家，牛價三千七百五十。

\cnshu 術曰：置七家共出一百九十、不足三百三十，九家共出二百七十、盈三十。令每家出率各以家數除之：一百九十除以七得二十七又七分之一；二百七十除以九得三十。乃以盈不足維乘之，并為實，并盈不足為法，實如法而一。

\cnzhu 此題以家為率，不以人為率。故先求每家所出之率，然後用盈不足術。七家出一百九十，則不足三百三十，是每家所出不足也；九家出二百七十，盈三十，是每家所出盈也。先通為每家之率，然後可用盈不足之法。此術之變，以家為率也。

\switchcolumn[1]
\enwen Now there is a group buying oxen. Seven households together contribute 190, falling short by 330; nine households together contribute 270, with an excess of 30. Question: how many households are there, and what is the price of the oxen?

\enda There are 126 households; the oxen cost 3750.

\enshu Method: Place 7 households contributing 190 with deficit 330, and 9 households contributing 270 with excess 30. Divide each total by the number of households: 190/7 gives 27 1/7; 270/9 gives 30. Then apply excess-deficit cross-multiplication: sum as dividend, sum of excess and deficit as divisor; divide.

\enzhu In this problem, households serve as the rate, not individual persons. Therefore first find the contribution rate per household, then apply the excess-deficit method. This is a variant of the method using households as rate.
\end{paracol}

\bigskip\hrule\bigskip

\subsection{第五題 / Problem 5: 米麥互換 / Rice and Wheat Exchange}

\begin{paracol}{2}
\switchcolumn[0]
\cnwen 今有米一斗，麥一斗，各值幾何？但知米三斗值麥二斗，麥四斗值米一斗。問：米、麥各一斗價若干？

\cnda 米一斗值四錢，麥一斗值一錢半。

\cnshu 術曰：以盈不足術求之。假令米一斗值四錢，則麥當值若干，驗其盈不足，乃以術求之。

\cnzhu 此術非直盈不足，乃以價為率，互相推求。假令一值，驗其盈肭，然後設為兩端，用盈不足術以齊之。米三斗值麥二斗，則米一斗值麥三分之二斗。麥四斗值米一斗，則麥一斗值米四分之一斗。以盈不足術通之，設米一斗值若干，驗其與麥價之盈肭，維乘求之。此以價為隱，而以盈不足顯之也。

\switchcolumn[1]
\enwen Now there is one dou of rice and one dou of wheat. What is each worth? It is known that 3 dou of rice is worth 2 dou of wheat, and 4 dou of wheat is worth 1 dou of rice. Question: what is the price of 1 dou of rice and 1 dou of wheat?

\enda One dou of rice is worth 4 cash; one dou of wheat is worth 1.5 cash.

\enshu Method: Apply the excess-deficit method. Assume one dou of rice is worth 4 cash, then determine what wheat should be worth, check the excess or deficit, and solve by the method.

\enzhu This method is not purely excess-deficit; it uses prices as rates and mutually deduces them. Assume one price, check its excess or deficit against the other, then set up two ends and use the excess-deficit method to harmonize them. This makes the hidden price relationship visible through excess and deficit.
\end{paracol}

\bigskip\hrule\bigskip

\subsection{第六題 / Problem 6: 良駑與驢 / Fine Horse, Nag, and Donkey}

\begin{paracol}{2}
\switchcolumn[0]
\cnwen 今有良駑二馬，與驢俱牽。良馬與駑馬共引之，良馬日引四十里，駑馬日引四十里。驢在後逐之，不及。問：驢日行幾何里？

\cnda 驢日行二十里。

\cnshu 術曰：以盈不足術求之。假令驢行三十里，則驢所行多於良駑之半，是為盈；假令驢行十里，則不及半，是為不足。維乘并實，如法求之。

\cnzhu 良駑俱行四十里，其半二十里。驢逐其後，當與之半相當。假令三十里，則多十里，為盈；假令十里，則少十里，為不足。故以盈不足術求之，得二十里。此術之意：良駑俱行四十里，驢逐其半，故以二十里為中平之數。

\switchcolumn[1]
\enwen Now there are a fine horse and a nag, both pulling together with a donkey chasing behind. The fine horse and nag each pull 40 li per day. The donkey follows behind but cannot catch up. Question: how many li does the donkey travel per day?

\enda The donkey travels 20 li per day.

\enshu Method: Apply the excess-deficit method. Assume the donkey travels 30 li: then it exceeds half of the horses' distance, which is excess. Assume it travels 10 li: then it falls short of half, which is deficit. Cross-multiply, sum as dividend, and solve by the method.

\enzhu The fine horse and nag both travel 40 li; half of this is 20 li. The donkey chases behind, and should match this half. Assuming 30 li gives 10 li more, which is excess; assuming 10 li gives 10 li less, which is deficit. Therefore apply the excess-deficit method to obtain 20 li.
\end{paracol}

\bigskip\hrule\bigskip

\subsection{第七題 / Problem 7: 垣高術 / Wall Height}

\begin{paracol}{2}
\switchcolumn[0]
\cnwen 今有垣不知高。立一表長五尺，去垣五丈，目距地五尺，望垣頂與表端參合。問：垣高幾何？

\cnda 垣高五丈五尺。

\cnshu 術曰：以盈不足術求之。假令垣高五丈，則視線當過表端若干；假令垣高六丈，則視線不及表端若干。維乘求之。

\cnzhu 此術本屬勾股，今以盈不足術驗之。去垣五丈，表高五尺，目高五尺，則表與目平。望垣頂與表端參合，是為表高與垣高之比例，如去垣之遠與目至垣頂之比例也。以相似勾股求之即得。此術雖屬勾股，而以盈不足驗之，所以明術之相通也。

\switchcolumn[1]
\enwen Now there is a wall of unknown height. A pole 5 chi tall is erected 5 zhang from the wall. The eye is 5 chi from the ground. Looking at the top of the wall aligns with the top of the pole. Question: how high is the wall?

\enda The wall is 5 zhang 5 chi high.

\enshu Method: Apply the excess-deficit method. Assume the wall is 5 zhang: then the sight line passes above the pole. Assume the wall is 6 zhang: then the sight line falls below the pole. Cross-multiply and solve.

\enzhu This method fundamentally belongs to right triangles (Gougu); here it is verified by the excess-deficit method. The alignment of the wall top with the pole top forms a proportion. Using similar right triangles yields the answer. Although this belongs to right triangles, using excess-deficit to verify it demonstrates the interconnectedness of methods.
\end{paracol}

\bigskip\hrule\bigskip

\subsection{第八題 / Problem 8: 五家共井 / Five Families Sharing a Well}

\begin{paracol}{2}
\switchcolumn[0]
\cnwen 今有五家共井。甲二綆不足，如乙一綆；乙三綆不足，如丙一綆；丙四綆不足，如丁一綆；丁五綆不足，如戊一綆；戊六綆不足，如甲一綆。問：井深、綆長各幾何？

\cnda 井深七丈二尺一寸。甲綆長二丈六尺五寸，乙綆長一丈九尺一寸，丙綆長一丈四尺八寸，丁綆長一丈二尺九寸，戊綆長七尺六寸。

\cnshu 術曰：此以盈不足為問，實則以方程入之。各以其綆之數，不足之數，列為方程，正負相生，求之即得。

\cnzhu 此題本以盈不足為問，而實非盈不足之術所能獨盡也。五家之綆長短不同，各家取綆若干而不足若干，以補他家一綆之數。此五綆之長與井深，共為六未知數，而題中僅給五條件，故為不定問題也。其答云者，特其一組正整數解耳。若以方程術解之，當令井深為一，求五家綆長之比，然後以井深定其實長。此題之精妙，在於以盈不足為表，以方程為裡，表裡相濟，乃得其真。

\switchcolumn[1]
\enwen Now five families share a well. Two lengths of A's rope are insufficient, and the shortfall equals one length of B's rope; three lengths of B's are insufficient, equaling one length of C's; four lengths of C's are insufficient, equaling one length of D's; five lengths of D's are insufficient, equaling one length of E's; six lengths of E's are insufficient, equaling one length of A's. Question: what is the depth of the well, and what is the length of each family's rope?

\enda The well is 7 zhang 2 chi 1 cun deep. A's rope is 2 zhang 6 chi 5 cun; B's is 1 zhang 9 chi 1 cun; C's is 1 zhang 4 chi 8 cun; D's is 1 zhang 2 chi 9 cun; E's is 7 chi 6 cun.

\enshu Method: Although posed as excess-deficit, this problem is actually solved by the rectangular arrays (Fangcheng) method. List each rope quantity and deficit quantity as equations; let positive and negative interact; solving yields the answer.

\enzhu This problem is phrased in terms of excess and deficit, but the excess-deficit method alone cannot fully solve it. The five rope lengths plus the well depth make six unknowns, while the problem only provides five conditions; hence it is an indeterminate problem. The subtlety lies in using excess-deficit as the exterior and Fangcheng as the interior; exterior and interior working together obtain the truth.
\end{paracol}

\bigskip\hrule\bigskip

\subsection{第九題 / Problem 9: 金與銀 / Gold and Silver}

\begin{paracol}{2}
\switchcolumn[0]
\cnwen 今有金九枚，銀十一枚，稱之適等。交易其一，金輕十三兩。問：金、銀一枚各重幾何？

\cnda 金重二斤三兩一十八銖，銀重一斤十三兩六銖。

\cnshu 術曰：以盈不足術求之。假令金重三斤，則九金二十七斤；令銀與之等，則每銀當重二十七斤十一分。交易其一，驗其輕重，得盈不足之數。再設一率，維乘求之。

\cnzhu 金九銀十一稱之適等者，總重等也。交易其一者，以一金換一銀也。金輕十三兩者，換後金方輕於原銀方十三兩也。以方程術解之尤便：設金重$x$，銀重$y$，則$9x = 11y$，且$8x + y = 10y + x - 13$兩。以盈不足術者，兩設其重，驗其差也。

\switchcolumn[1]
\enwen Now there are 9 pieces of gold and 11 pieces of silver; they balance exactly. Exchange one piece, and the gold side is lighter by 13 liang. Question: how much does one piece of gold and one piece of silver weigh?

\enda One piece of gold weighs 2 jin 3 liang 18 zhu; one piece of silver weighs 1 jin 13 liang 6 zhu.

\enshu Method: Apply the excess-deficit method. Assume gold weighs 3 jin: then 9 pieces weigh 27 jin; making silver balance this, each piece should weigh 27/11 jin. Exchange one piece and check the difference to obtain the excess-deficit numbers. Make a second assumption, cross-multiply and solve.

\enzhu Nine pieces of gold and eleven pieces of silver balancing means their total weights are equal. Exchanging one piece means swapping one gold for one silver. Gold being 13 liang lighter means after exchange, the gold side is lighter than the original silver side by 13 liang. This is especially conveniently solved by the Fangcheng method.
\end{paracol}

\bigskip\hrule\bigskip

\subsection{第十題 / Problem 10: 程傳委輸 / Transporting Grain}

\begin{paracol}{2}
\switchcolumn[0]
\cnwen 今有程傳委輸，空車日行七十里，重車日行五十里。今載太倉粟輸上林，五日三返。問：太倉去上林幾何？

\cnda 太倉去上林四十八里一百八十分里之十一。

\cnshu 術曰：以盈不足術求之。假令去四十五里，則五日三返不足若干里；假令去五十里，則五日三返盈若干里。維乘并實，如法求之。

\cnzhu 空車日行七十里者，往而不載也；重車日行五十里者，載粟而返也。五日三返者，一往一返為一返，三返則三往三返也。此以盈不足術驗之者，取其近似也。此術雖以盈不足入之，實則以調和平均為本也。

\switchcolumn[1]
\enwen Now there is a relay transport task. An empty cart travels 70 li per day; a loaded cart travels 50 li per day. Grain is to be transported from Taicang to Shanglin, with three round trips in five days. Question: what is the distance from Taicang to Shanglin?

\enda The distance is 48 11/180 li.

\enshu Method: Apply the excess-deficit method. Assume the distance is 45 li: then three round trips in five days falls short. Assume 50 li: then three round trips exceeds. Cross-multiply, sum as dividend, and solve.

\enzhu An empty cart traveling 70 li per day is going without load; a loaded cart traveling 50 li per day is returning with grain. Three round trips in five days means one round trip is one going and one returning. Although this uses the excess-deficit method, its foundation is actually the harmonic mean.
\end{paracol}

\bigskip\hrule\bigskip

\subsection{第十一題 / Problem 11: 鹿與兔 / Deer and Rabbit}

\begin{paracol}{2}
\switchcolumn[0]
\cnwen 今有鹿直西走，兔直東走。兔行一百步，鹿行七十步。兔在鹿東一百五十步，同時發。問：幾何步相逢？

\cnda 兔行二百五十步，鹿行一百七十五步，相逢。

\cnshu 術曰：以盈不足術求之。假令兔行二百步，則鹿行一百四十步，驗其盈不足之數。再假令一行數，維乘求之。

\cnzhu 兔東鹿西，相向而行，故為相逢之術。兔行一百步，鹿行七十步，則兔鹿步率之比為十比七也。相距一百五十步，同時發，則相逢之時，兔行十份，鹿行七份，共十七份。以一百五十步分為十七份，每份150/17步。

\switchcolumn[1]
\enwen Now a deer runs straight west and a rabbit runs straight east. The rabbit runs 100 paces while the deer runs 70 paces. The rabbit is 150 paces east of the deer; they start at the same time. Question: after how many paces do they meet?

\enda The rabbit runs 250 paces, the deer runs 175 paces, and they meet.

\enshu Method: Apply the excess-deficit method. Assume the rabbit runs 200 paces: then the deer runs 140 paces, and check the excess-deficit numbers. Make a second assumption, cross-multiply and solve.

\enzhu The rabbit runs east and the deer runs west; they move toward each other. The ratio of their pace rates is 100:70 = 10:7. Separated by 150 paces and starting simultaneously, at meeting time the rabbit covers 10 parts and the deer 7 parts, totaling 17 parts.
\end{paracol}

\bigskip\hrule\bigskip

\subsection{第十二題 / Problem 12: 蒲與莞 / Rush and Cattail}

\begin{paracol}{2}
\switchcolumn[0]
\cnwen 今有蒲生一日，長三尺；莞生一日，長一尺。蒲生日自半，莞生日自倍。問：幾何日而長等？

\cnda 二日十三分日之六，而長等。各長四尺五寸十三分寸之五。

\cnshu 術曰：以盈不足術求之。假令二日，不足一尺五寸；令之三日，盈一尺七寸半。維乘并實，如法求之。

\cnzhu 蒲日生三尺而自半者，第一日長三尺，第二日長一尺五寸，第三日長七寸半，是為日自半也。莞日生一尺而自倍者，第一日長一尺，第二日長二尺，第三日長四尺，是為日自倍也。此術之巧，在於兩物生長之率不同，一自半而一自倍，以盈不足術齊而同之。

\switchcolumn[1]
\enwen Now a rush grows 3 chi on the first day; a cattail grows 1 chi on the first day. The rush's growth halves each day; the cattail's growth doubles each day. Question: after how many days are their lengths equal?

\enda After 2 6/13 days their lengths are equal. Each is 4 5/13 chi.

\enshu Method: Apply the excess-deficit method. Assume 2 days: deficit of 1 chi 5 cun. Assume 3 days: excess of 1 chi 7.5 cun. Cross-multiply, sum as dividend, and solve.

\enzhu The rush growing 3 chi per day and halving means: first day 3 chi, second day 1.5 chi, third day 7.5 cun---this is daily halving. The cattail growing 1 chi per day and doubling means: first day 1 chi, second day 2 chi, third day 4 chi---this is daily doubling. The ingenuity lies in the two objects having different growth rates, and the excess-deficit method equalizes them.
\end{paracol}

\bigskip\hrule\bigskip

\subsection{第十五題 / Problem 15: 漆三油四 / Lacquer and Oil}

\begin{paracol}{2}
\switchcolumn[0]
\cnwen 今有漆三得油四，油四和漆五。今有漆三斗，欲令分以易油，還和餘漆。問：出漆、得油、和漆各幾何？

\cnda 出漆一斗一升四分升之一，得油一斗五升，和漆一斗八升四分升之三。

\cnshu 術曰：以盈不足術求之。假令出漆一斗，不足若干；令出漆一斗二升，盈若干。維乘并實，如法求之。

\cnzhu 漆三得油四者，以漆三斗易油四斗也。油四和漆五者，以油四斗和漆五斗也。今有漆三斗，分其漆以易油，又以所得之油和其餘漆。此術以漆易油，以油和漆，往還相求，以盈不足定其率。

\switchcolumn[1]
\enwen Now 3 parts of lacquer yield 4 parts of oil; 4 parts of oil mix with 5 parts of lacquer. Now there are 3 dou of lacquer. It is desired to divide some lacquer to exchange for oil, then use the obtained oil to mix with the remaining lacquer. Question: how much lacquer is used, how much oil is obtained, and how much lacquer-oil mixture results?

\enda Lacquer used: 1 dou 1 sheng 2.25 sheng; oil obtained: 1 dou 5 sheng; mixture: 1 dou 8 sheng 3.75 sheng.

\enshu Method: Apply the excess-deficit method. Assume 1 dou of lacquer is used: deficit. Assume 1 dou 2 sheng: excess. Cross-multiply, sum as dividend, and solve.

\enzhu This method exchanges lacquer for oil and uses oil to mix lacquer, going back and forth, using excess-deficit to fix the rate, demonstrating the inexhaustible versatility of the method.
\end{paracol}

\bigskip\hrule\bigskip

\subsection{第十六題 / Problem 16: 惡粟為米 / Coarse Millet into Rice}

\begin{paracol}{2}
\switchcolumn[0]
\cnwen 今有惡粟二十斗，春之為米，米率五斗。今欲為糲米、粺米各幾何？

\cnda 糲米十斗，粺米六斗六升三分升之二。

\cnshu 術曰：以盈不足術求之。假令糲米十斗，則粺米當若干，驗其盈不足。再設一率，維乘求之。

\cnzhu 惡粟二十斗，米率五斗者，一斗惡粟春得五升米也。糲米率五，粺米率六。以盈不足術者，兩設糲米之斗數，求粺米之盈肭也。此以米率之隱，而以盈不足顯之也。

\switchcolumn[1]
\enwen Now there are 20 dou of coarse millet; pounding it yields rice at a rate of 5 sheng per dou. It is desired to divide this into roughly-husked rice (li mi) and finely-husked rice (bai mi). Question: how much of each?

\enda Roughly-husked rice: 10 dou; finely-husked rice: 6 dou 6 sheng 2/3 sheng.

\enshu Method: Apply the excess-deficit method. Assume 10 dou of roughly-husked rice: then determine the finely-husked rice and check the excess or deficit. Make a second assumption, cross-multiply and solve.

\enzhu This reveals the hidden rice rates through the excess-deficit method.
\end{paracol}

\bigskip\hrule\bigskip

\subsection{第十七題 / Problem 17: 持米出關 / Carrying Rice Through Passes}

\begin{paracol}{2}
\switchcolumn[0]
\cnwen 今有人持米出三關，外關三而取一，中關五而取一，內關七而取一，餘米五斗。問：本持米幾何？

\cnda 本持米十斗九升八分升之三。

\cnshu 術曰：以盈不足術求之。假令持米八斗，不足若干；令持米十二斗，盈若干。維乘并實，如法求之。

\cnzhu 三關各取其幾分之一。外關三取一，是取其三分之一也；中關五取一，是取其五分之一也；內關七取一，是取其七分之一也。過三關之後，餘米五斗。此術以三關之稅率為隱，以盈不足顯之，所以見術之無幽不顯也。

\switchcolumn[1]
\enwen Now a person carries rice through three passes. The outer pass takes one-third; the middle pass takes one-fifth; the inner pass takes one-seventh. After passing through all three, 5 dou of rice remain. Question: how much rice did the person originally carry?

\enda The original amount was 10 dou 9 sheng 3/8 sheng.

\enshu Method: Apply the excess-deficit method. Assume 8 dou were carried: deficit. Assume 12 dou: excess. Cross-multiply, sum as dividend, and solve.

\enzhu Each pass takes a certain fraction. This method makes the hidden tax rates of the three passes visible through excess and deficit, demonstrating that the method reveals all that is hidden.
\end{paracol}

\bigskip\hrule\bigskip

\subsection{第十八題 / Problem 18: 五人分五錢 / Five People Dividing Five Cash}

\begin{paracol}{2}
\switchcolumn[0]
\cnwen 今有五人分五錢，令上二人所得與下三人等。問：各得幾何？

\cnda 甲得一錢六分钱之二，乙得一錢六分钱之一，丙得一錢，丁得六分钱之五，戊得六分钱之四。

\cnshu 術曰：以盈不足術求之。假令五人各得一錢，則上二人得二錢，下三人得三錢，盈一錢。此當以衰分術入之。

\cnzhu 五人分五錢而令上二人與下三人等者，上二人所得當為二錢半，下三人所得亦當為二錢半也。然此題以衰分術為正：令五人所得為差率，以差分配之。此術雖以盈不足入之，實以衰分為本也。

\switchcolumn[1]
\enwen Now five people divide 5 cash, such that the top two people's share equals the bottom three people's share. Question: how much does each receive?

\enda A receives 1 2/6 cash; B receives 1 1/6 cash; C receives 1 cash; D receives 5/6 cash; E receives 4/6 cash.

\enshu Method: Apply the excess-deficit method. Assume each receives 1 cash: then the top two get 2 cash and the bottom three get 3 cash, excess of 1 cash. [Note: properly solved by proportional division.]

\enzhu This problem is properly solved by the proportional division method. Although this uses the excess-deficit method, its foundation is actually proportional division.
\end{paracol}

\bigskip\hrule\bigskip

\subsection{第十九題 / Problem 19: 弧田徑術 / Arc-Field Area}

\begin{paracol}{2}
\switchcolumn[0]
\cnwen 今有弧田，弦三十步，矢十五步。問：為田幾何？

\cnda 一畝九十七步半。

\cnshu 術曰：以弦乘矢，矢又自乘，并之，三而一。

\cnzhu 弧田術者，以弦矢求弧田之面積也。弦乘矢者，以弦為廣，矢為縱，相乘得長方形之積。矢自乘者，以矢為方之邊，自乘得正方之積。并之者，合兩積為一。三而一者，以弧田之積約當此兩積和之三分之一也。此術為弧田之近似術。

\switchcolumn[1]
\enwen Now there is an arc-field with chord 30 paces and arrow (sagitta) 15 paces. Question: what is the area?

\enda 1 mu 97.5 square paces.

\enshu Method: Multiply the chord by the arrow; also square the arrow; sum these; divide by 3.

\enzhu This is an approximation method for arc-fields. The arc-field is the region enclosed by a circular arc and its chord; its shape is curved and convex, not fully expressible by straight lines.
\end{paracol}

\bigskip\hrule\bigskip

\subsection{第二十題 / Problem 20: 環田徑術 / Annulus Field}

\begin{paracol}{2}
\switchcolumn[0]
\cnwen 今有環田，中周六十二步，外周四步，徑十二步。問：為田幾何？

\cnda 二畝五十五步。

\cnshu 術曰：并中外周而半之，以徑乘之為積。

\cnzhu 環田者，猶圓田之中去一圓也。并中外周而半之者，取環之平均周長也。以徑乘之者，以平均周長為廣，徑為縱也。此術以環田近似為矩形，故以廣縱相乘得積。若精密求之，當以兩圓面積之差求之。

\switchcolumn[1]
\enwen Now there is an annulus field with inner circumference 62 paces, outer circumference [corrected], and diameter [thickness] 12 paces. Question: what is the area?

\enda 2 mu 55 square paces.

\enshu Method: Sum the inner and outer circumferences and halve; multiply by the diameter to get the area.

\enzhu This method approximates the annulus as a rectangle. For a precise calculation, one should use the difference of the two circle areas.
\end{paracol}

\bigskip\hrule\bigskip

\subsection{第二十一題 / Problem 21: 三人共車 / Three People Per Cart}

\begin{paracol}{2}
\switchcolumn[0]
\cnwen 今有三人共車，二車空；二人共車，九人步。問：車幾何？人幾何？

\cnda 車一十五乘，人三十九人。

\cnshu 術曰：以盈不足術求之。假令車十乘，則三人共車載三十人，餘九人步，是為盈九。假令車二十乘，則二人共車載四十人，缺一人，是為不足一。維乘并實，如法求之。

\cnzhu 三人共車二車空者，每車三人則二車無人乘也。二人共車九人步者，每車二人則餘九人不行也。此術以車之虛實為盈肭，亦變術之巧者也。

\switchcolumn[1]
\enwen Now if three people share each cart, two carts are empty; if two people share each cart, nine people must walk. Question: how many carts and how many people?

\enda There are 15 carts and 39 people.

\enshu Method: Apply the excess-deficit method. Assume 10 carts: then three per cart carries 30 people, leaving 9 walking---excess of 9. Assume 20 carts: then two per cart carries 40 people, one short---deficit of 1. Cross-multiply, sum as dividend, and solve.

\enzhu This method uses the emptiness or fullness of carts as excess and deficit---a clever variation of the method.
\end{paracol}

\bigskip\hrule\bigskip

\subsection{第二十二題 / Problem 22: 綿布互易 / Cotton and Cloth Exchange}

\begin{paracol}{2}
\switchcolumn[0]
\cnwen 今有綿一斤直錢二百四十，布一匹直錢三百二十。今有綿五斤，欲易布幾何？

\cnda 布三匹七分匹之五。

\cnshu 術曰：以盈不足術求之。假令易三匹，則綿價一千二百，布價九百六十，盈二百四十。假令易四匹，則綿價一千二百，布價一千二百八十，不足八十。維乘并實，如法求之。

\cnzhu 綿一斤直二百四十，五斤則直一千二百。布一匹直三百二十。以盈不足術者，兩設易布之匹數，驗其綿布價值之盈肭也。

\switchcolumn[1]
\enwen Now cotton costs 240 cash per jin; cloth costs 320 cash per bolt. With 5 jin of cotton, how much cloth can be obtained?

\enda 3 5/7 bolts.

\enshu Method: Apply the excess-deficit method. Assume exchanging for 3 bolts: cotton value 1200, cloth value 960, excess of 240. Assume 4 bolts: cotton value 1200, cloth value 1280, deficit of 80. Cross-multiply, sum as dividend, and solve.

\enzhu Using the excess-deficit method: make two assumptions about the number of bolts and check the excess or deficit of the cotton-cloth values.
\end{paracol}

\bigskip\hrule\bigskip

\subsection{第二十三題 / Problem 23: 粟米貴賤 / Millet and Rice Grades}

\begin{paracol}{2}
\switchcolumn[0]
\cnwen 今有粟一斗，欲為糲米、粺米、糞米三等。糲米率三十，粺米率二十七，糞米率二十四。問：各為米幾何？

\cnda 糲米三升，粺米二升七合，糞米二升四合。

\cnshu 術曰：以盈不足術入衰分求之。假令糲米三升，則粺米、糲米遞減之，驗其盈不足，維乘求之。

\cnzhu 糲米率三十者，粟一斗為糲米三升也。粺米率二十七者，為粺米二升七合也。糞米率二十四者，為糞米二升四合也。此術以三米之率為隱，以盈不足為顯。

\switchcolumn[1]
\enwen Now there is 1 dou of millet, to be divided into three grades of rice: roughly-husked (li), finely-husked (bai), and polished (fen). Rates: li 30, bai 27, fen 24. Question: how much of each?

\enda Li mi: 3 sheng; bai mi: 2 sheng 7 ge; fen mi: 2 sheng 4 ge.

\enshu Method: Apply the excess-deficit method within proportional division. Assume 3 sheng of li mi: then bai and fen mi decrease accordingly; check the excess or deficit; cross-multiply and solve.

\enzhu This method makes the three rice rates visible through excess and deficit.
\end{paracol}

\bigskip\hrule\bigskip

\subsection{第二十四題 / Problem 24: 土功程效 / Earthwork Labor}

\begin{paracol}{2}
\switchcolumn[0]
\cnwen 今有築堤，甲縣人功一日築七尺，乙縣人功一日築五尺。今令兩縣共築，甲縣先築一日，乙縣繼之。問：幾何日而成？

\cnda 甲築一日，乙築二日，共成。

\cnshu 術曰：以盈不足術求之。假令共築二日，則甲築十四尺，乙築五尺，共十九尺，驗其盈不足。再設一日，維乘求之。

\cnzhu 甲縣人功一日七尺，乙縣人功一日五尺。甲先築一日，築七尺。餘令乙築，乙一日五尺，則餘二尺。甲復築一日，又七尺，共十九尺。此術雖以盈不足入之，實以衰分為本也。

\switchcolumn[1]
\enwen Now a dike is to be built. County A's workers build 7 chi per day; County B's build 5 chi per day. Both counties build together; County A works first for one day, then County B continues. Question: how many days until completion?

\enda County A builds for 1 day; County B builds for 2 days; together they complete it.

\enshu Method: Apply the excess-deficit method. Assume both build for 2 days total: then A builds 14 chi and B builds 5 chi, total 19 chi; check the excess or deficit. Make a second assumption, cross-multiply and solve.

\enzhu County A's laborers build 7 chi per day; County B's build 5 chi per day. Although this uses the excess-deficit method, its foundation is actually proportional division.
\end{paracol}

\bigskip\hrule\bigskip

%% ========================================================================
%% VOLUME 8: FANGCHENG (方程)
%% ========================================================================
\newpage
\section*{卷第八\quad 方程\quad /\quad Volume 8\quad Rectangular Arrays}
\addcontentsline{toc}{section}{卷第八 方程 / Volume 8 Rectangular Arrays}

\begin{paracol}{2}
\switchcolumn[0]
\begin{quote}\color{quotecolor}\small
\textit{方程章乃中國古代代數之巔峰成就。其提出解線性方程組之系統方法——本質即高斯消元法——較高斯早一千五百餘年。本章亦包含世界數學史上負數之最早系統運用。}
\end{quote}
\switchcolumn[1]
\begin{quote}\color{quotecolor}\small
\textit{The Fangcheng chapter is the crowning achievement of ancient Chinese algebra. It presents a systematic method for solving systems of linear equations --- essentially Gaussian elimination --- over 1500 years before Gauss. This chapter also contains the first known use of negative numbers in world mathematics.}
\end{quote}
\end{paracol}

\subsection*{General Method / 方程術總論}

\begin{paracol}{2}
\switchcolumn[0]
\cnshu 方程術曰：置上禾三秉，中禾二秉，下禾一秉，實三十九；上禾二秉，中禾三秉，下禾一秉，實三十四；上禾一秉，中禾二秉，下禾三秉，實二十六。以右行上禾遍乘中行，而以直除。又乘其次，亦以直除。然以中行中禾不盡者遍乘左行，而以直除。左方下禾不盡者，上為法，下為實。實即下禾之實。

\cnzhu 程，課程也。群物總雜，各列有數，總言其實。令每行為率，二物者再程，三物者三程，皆如物數程之，並列為行，故謂之方程。行之左右，無所同存，且為有所據而言耳。

\switchcolumn[1]
\enshu The Method of Rectangular Arrays says: Place [coefficients] in rows. Multiply the pivot row throughout and subtract directly from rows below. Repeat for each column. The remaining coefficient is the divisor, the shi is the dividend. Dividing gives the unknown's value.

\enzhu ``Cheng'' means course or standard. When various items are mixed together, each is listed with its quantity and the total is given as ``shi.'' Each row forms a rate; as many courses as there are items---arranged together as rows; hence ``rectangular arrays.''
\end{paracol}

\begin{mathdisplay}
\textbf{Key Steps:} 1) Arrange coefficients in rows; 2) Multiply pivot row and subtract from rows below (直除); 3) Repeat for each column; 4) Back-substitute.
\textbf{Historical Note:} Liu Hui used red rods for positive numbers (正) and black rods for negative numbers (負).
\end{mathdisplay}

\bigskip\hrule\bigskip

\subsection{第一題 / Problem 1: 三禾求實 / Three Grains' Yield}

\begin{paracol}{2}
\switchcolumn[0]
\cnwen 今有上禾三秉，中禾二秉，下禾一秉，實三十九斗；上禾二秉，中禾三秉，下禾一秉，實三十四斗；上禾一秉，中禾二秉，下禾三秉，實二十六斗。問：上、中、下禾實一秉各幾何？

\cnda 上禾一秉，九斗四分斗之一；中禾一秉，四斗四分斗之一；下禾一秉，二斗四分斗之三。

\cnshu 術曰：方程術。置上禾、中禾、下禾及實於右、中、左三行。以右行上禾三遍乘中行，以直除之，消去中行上禾。又以右行上禾三遍乘左行，以直除之，消去左行上禾。然後以中行中禾不盡者遍乘左行，以直除之，消去左行中禾。

\cnzhu 此方程術之正也。置三行於位：右行上禾三、中禾二、下禾一、實三十九；中行上禾二、中禾三、下禾一、實三十四；左行上禾一、中禾二、下禾三、實二十六。以右行上禾三乘中行，得六、九、三、一百零二；與右行直除之。

\switchcolumn[1]
\enwen Now there are 3 bundles of upper grain, 2 of middle, 1 of lower, yielding 39 dou; 2 of upper, 3 of middle, 1 of lower yielding 34 dou; 1 of upper, 2 of middle, 3 of lower yielding 26 dou. Question: how much does one bundle of each grain yield?

\enda Upper: 9 1/4 dou; Middle: 4 1/4 dou; Lower: 2 3/4 dou.

\enshu Method: Fangcheng method. Arrange upper, middle, lower grain and shi in right, middle, left rows. Multiply middle row by upper grain coefficient 3 of right row; subtract directly to eliminate upper grain. Repeat for left row. Then eliminate middle grain from left row. Finally divide remaining lower grain coefficient into shi.

\enzhu This is the orthodox Fangcheng method. Arrange three rows: right (3 upper, 2 middle, 1 lower, shi 39); middle (2 upper, 3 middle, 1 lower, shi 34); left (1 upper, 2 middle, 3 lower, shi 26). Eliminate by direct subtraction.
\end{paracol}

\bigskip\hrule\bigskip

\subsection{第二題 / Problem 2: 五家共輸 / Five Families Contributing Grain}

\begin{paracol}{2}
\switchcolumn[0]
\cnwen 今有甲、乙、丙、丁、戊五家共輸粟。甲所輸加乙之一半，乙所輸加丙之三分之一，丙所輸加丁之四分之一，丁所輸加戊之五分之一，戊所輸加甲之六分之一，各適等。問：各家所輸幾何？

\cnda 各家所輸均為一斛。甲輸二十五分斛之十二，乙輸二十五分斛之四，丙輸二十五分斛之三，丁輸二十五分斛之二，戊輸二十五分斛之一。

\cnshu 術曰：以方程術入之。令各家所輸為未知數，以條件列方程，正負相生，求之即得。

\cnzhu 此題以方程術入之為便。令甲、乙、丙、丁、戊所輸分別為$a, b, c, d, e$。則有：$a + b/2 = b + c/3 = c + d/4 = d + e/5 = e + a/6$。令此等於$k$，則五式聯立，可以方程術解之。

\switchcolumn[1]
\enwen Now five families A, B, C, D, and E are to contribute millet. A's contribution plus half of B's; B's plus one-third of C's; C's plus one-fourth of D's; D's plus one-fifth of E's; E's plus one-sixth of A's---all equal. Question: how much does each family contribute?

\enda Each family's adjusted contribution equals 1 hu. The ratios are: A: 12/25, B: 4/25, C: 3/25, D: 2/25, E: 1/25.

\enshu Method: Apply the Fangcheng method. Let each family's contribution be unknowns; form equations from the conditions; let positive and negative interact; solving yields the answer.

\enzhu This problem is conveniently solved by the Fangcheng method. Let the contributions be $a, b, c, d, e$; then all expressions equal a constant $k$, giving five equations solvable by Fangcheng.
\end{paracol}

\bigskip\hrule\bigskip

\subsection{第三題 / Problem 3: 三畜直金 / Three Animals' Gold Value}

\begin{paracol}{2}
\switchcolumn[0]
\cnwen 今有牛二、羊五，直金十兩；牛三、豕三，直金九兩；羊六、豕八，直金八兩。問：牛、羊、豕各直金幾何？

\cnda 牛一直金一兩二十一分兩之十三，羊一直金二十一分兩之二十，豕一直金二十一分兩之四。

\cnshu 術曰：方程術。置牛、羊、豕之數及直金於三行，以直除消去牛、羊，得下豕之實，實如法得豕之直。以次求羊、牛之直。

\cnzhu 牛二、羊五直十兩，牛三、豕三直九兩，羊六、豕八直八兩。以方程術：置三行，先消去牛，次消去羊，得豕之實。以豕之實推羊，以羊之實推牛。此三物方程之正法也。

\switchcolumn[1]
\enwen Now 2 oxen and 5 sheep are worth 10 taels of gold; 3 oxen and 3 pigs are worth 9 taels; 6 sheep and 8 pigs are worth 8 taels. Question: how much is each animal worth?

\enda Ox: 1 13/21 taels; Sheep: 20/21 taels; Pig: 4/21 taels.

\enshu Method: Fangcheng method. Arrange oxen, sheep, pigs and gold values in three rows. Eliminate oxen and sheep by direct subtraction to obtain pig's shi; divide to get pig's value. Then find sheep and ox values.

\enzhu This is the orthodox Fangcheng method for three items: arrange three rows, first eliminate oxen, then sheep, obtaining pig's shi. From pig's shi deduce sheep's, and from sheep's deduce ox's.
\end{paracol}

\bigskip\hrule\bigskip

\subsection{第四題 / Problem 4: 麻麥稻菽黍 / Five Crops}

\begin{paracol}{2}
\switchcolumn[0]
\cnwen 今有麻九斗、麥七斗、菽三斗、答二斗、黍五斗，直錢一百四十；麻七斗、麥六斗、菽四斗、答五斗、黍三斗，直錢一百二十八；麻三斗、麥五斗、菽七斗、答六斗、黍四斗，直錢一百一十六；麻二斗、麥三斗、菽五斗、答七斗、黍六斗，直錢一百零四；麻五斗、麥四斗、菽六斗、答三斗、黍七斗，直錢一百一十二。問：五物各一斗直錢幾何？

\cnda 麻一斗七錢，麥一斗四錢，菽一斗三錢，答一斗五錢，黍一斗六錢。

\cnshu 術曰：方程術。以五物列五行，以直錢為實，遍乘直除，消去各物之率，各得其一斗之實。

\cnzhu 此五物方程，較之三物方程為繁。其術一也：列五行为程，以右行首物遍乘其下各行，以直除之，消去首物；次以中行次物遍乘其下各行，以直除之，消去次物；如是遞推，終得一物之實。

\switchcolumn[1]
\enwen Now [various quantities of] five crops---hemp, wheat, beans, malt, and millet---are worth certain amounts of cash. Question: how much cash is one dou of each crop worth?

\enda Hemp: 7 cash; Wheat: 4 cash; Beans: 3 cash; Malt: 5 cash; Millet: 6 cash per dou.

\enshu Method: Fangcheng method. Arrange five items in five rows with cash values as shi. Multiply and subtract directly to eliminate each item's coefficient.

\enzhu This five-item Fangcheng is more complex than the three-item version, but the method is the same: arrange five rows; eliminate leading item from rows below; proceed until one item's shi remains; back-substitute.
\end{paracol}

\bigskip\hrule\bigskip

\subsection{第五題 / Problem 5: 正負術 / Positive and Negative Numbers}

\begin{paracol}{2}
\switchcolumn[0]
\cnwen 今有上禾七秉，損實一斗，益之下禾二秉，則實滿十斗；下禾八秉，益實一斗，與上禾三秉，則實滿十斗。問：上、下禾一秉各幾何？

\cnda 上禾一秉一斗五十二分斗之四十七，下禾一秉五十二分斗之三十九。

\cnshu 術曰：以方程術入之。損實一斗者，正也；益實一斗者，負也。正負列位，如方程消之。

\cnzhu 正負術者，方程術之輔也。以兩算得失相反，要令正負以名之。正算赤，負算黑。否則以邪正為異。同名相除，異名相益。正無入負之，負無入正之。此正負之術也。

\switchcolumn[1]
\enwen Now 7 bundles of upper grain, with 1 dou removed, and adding 2 bundles of lower grain, gives exactly 10 dou; 8 bundles of lower grain, with 1 dou added, together with 3 bundles of upper grain, gives exactly 10 dou. Question: how much does one bundle of each grain yield?

\enda Upper grain: 1 47/52 dou; Lower grain: 39/52 dou.

\enshu Method: Apply the Fangcheng method. Removing 1 dou of shi is positive; adding 1 dou of shi is negative. Arrange positive and negative positions and eliminate.

\enzhu The positive-negative method is the auxiliary to the Fangcheng method. Positive counting rods are red, negative are black. Same names subtract; different names add. This is the first known systematic use of signed numbers in world mathematics.
\end{paracol}

\bigskip\hrule\bigskip

\subsection{第六題 / Problem 6: 四畜求直 / Horse, Ox, and Sheep}

\begin{paracol}{2}
\switchcolumn[0]
\cnwen 今有馬一、牛二、羊三，直金二十五兩；馬二、牛三、羊一，直金二十八兩；馬三、牛一、羊二，直金二十三兩。問：馬、牛、羊各直金幾何？

\cnda 馬一直金九兩，牛一直金五兩，羊一直金二兩。

\cnshu 術曰：方程術。列三行，馬、牛、羊之數及直金為實。以右行馬一遍乘中行、左行，以直除之，消去馬。次以中行牛遍乘左行，以直除之，消去牛。左行羊不盡者，上為法，下為實，實如法得羊之直。

\cnzhu 此三物方程之又一例也。馬一、牛二、羊三直二十五兩；馬二、牛三、羊一直二十八兩；馬三、牛一、羊二直二十三兩。以方程術消之：右行馬一為最簡，故以右行遍乘中行、左行，直除消馬。

\switchcolumn[1]
\enwen Now 1 horse, 2 oxen, and 3 sheep are worth 25 taels; 2 horses, 3 oxen, and 1 sheep worth 28 taels; 3 horses, 1 ox, and 2 sheep worth 23 taels. Question: how much is each animal worth?

\enda Horse: 9 taels; Ox: 5 taels; Sheep: 2 taels.

\enshu Method: Fangcheng method. Arrange three rows with quantities and gold values. Multiply middle and left rows by right row's horse coefficient 1; subtract directly to eliminate horses. Then eliminate oxen from left row.

\enzhu This is another three-item Fangcheng. The right row has horse coefficient 1, the simplest, so multiply and subtract to eliminate horses. Then eliminate oxen to obtain sheep's shi. Back-substitute to find oxen and horses.
\end{paracol}

\bigskip\hrule\bigskip

\subsection{第七題 / Problem 7: 粟麥互易 / Grain Exchange}

\begin{paracol}{2}
\switchcolumn[0]
\cnwen 今有上禾六秉，損實一斗八升，當下禾一十秉；下禾一十五秉，損實五升，當上禾五秉。問：上、下禾一秉各幾何？

\cnda 上禾一秉八升，下禾一秉三升。

\cnshu 術曰：方程術。列上禾、下禾之數及損實為行。以正負術記之，直除消去，求得各秉之實。

\cnzhu 上禾六秉損一斗八升當下禾十秉者，六秉上禾之實減一斗八升，等於十秉下禾之實也。下禾十五秉損五升當上禾五秉者，十五秉下禾之實減五升，等於五秉上禾之實也。此二物二程之方程，較之三物為簡，然其法不異。

\switchcolumn[1]
\enwen Now 6 bundles of upper grain, with 1 dou 8 sheng removed, equals 10 bundles of lower grain; 15 bundles of lower grain, with 5 sheng removed, equals 5 bundles of upper grain. Question: how much does one bundle of each grain yield?

\enda Upper grain: 8 sheng; Lower grain: 3 sheng per bundle.

\enshu Method: Fangcheng method. Arrange upper grain, lower grain quantities and removed shi as rows. Record with positive-negative; eliminate by direct subtraction; find yield per bundle.

\enzhu This two-item, two-equation Fangcheng is simpler than the three-item case, but the method does not differ.
\end{paracol}

\bigskip\hrule\bigskip

\subsection{第八題 / Problem 8: 五家共輸稅 / Five Classes Contributing Tax}

\begin{paracol}{2}
\switchcolumn[0]
\cnwen 今有五等之家共輸稅。甲五家、乙四家、丙三家、丁二家、戊一家，共輸稅一千斛。欲令各依等衰輸之。問：各輸幾何？

\cnda 甲輸二百七十七斛七分斛之六，乙輸二百二十二斛七分斛之二，丙輸一百六十六斛七分斛之四，丁輸一百一十一斛七分斛之一，戊輸五十五斛七分斛之五。

\cnshu 術曰：以方程術入之。令五家為五未知，以家數及共輸為實，依衰分列方程，消去求之。

\cnzhu 五等之家，各有多少之異。甲五、乙四、丙三、丁二、戊一，共十五。以千斛分為十五分，甲得其五，乙得其四，丙得其三，丁得其二，戊得其一。此術雖以方程入之，實以衰分為本也。

\switchcolumn[1]
\enwen Now five classes of households are to contribute tax: A has 5 households, B has 4, C has 3, D has 2, E has 1; together they contribute 1000 hu. It is desired that each class contribute proportionally. Question: how much does each class contribute?

\enda Class A: 277 6/7 hu; B: 222 2/7 hu; C: 166 4/7 hu; D: 111 1/7 hu; E: 55 5/7 hu.

\enshu Method: Apply the Fangcheng method. Let the five classes be five unknowns; use household counts and total contribution as shi; form equations according to proportional decrease and eliminate.

\enzhu Five classes differ in quantity. The proportions A:B:C:D:E = 5:4:3:2:1 sum to 15. Divide 1000 hu accordingly. Although this uses the Fangcheng method, its foundation is proportional division.
\end{paracol}

\bigskip\hrule\bigskip

\subsection{第九題 / Problem 9: 三人持錢 / Three People Holding Money}

\begin{paracol}{2}
\switchcolumn[0]
\cnwen 今有甲、乙、丙三人持錢。甲得乙之半而錢滿四十八，乙得丙之半而錢滿四十八，丙得甲之半而錢滿四十八。問：甲、乙、丙各持錢幾何？

\cnda 甲持錢三十六，乙持錢二十四，丙持錢十二。

\cnshu 術曰：以方程術入之。令甲、乙、丙持錢為三未知數，以條件列三方程，消去求之。

\cnzhu 甲得乙之半而滿四十八者，甲加乙之半等於四十八也。設甲持$x$，乙持$y$，丙持$z$，則：$x + y/2 = 48$，$y + z/2 = 48$，$z + x/2 = 48$。以方程術列三行，正負記之，消去求之。

\switchcolumn[1]
\enwen Now three people A, B, and C hold money. If A obtains half of B's money, he has exactly 48; if B obtains half of C's money, he has exactly 48; if C obtains half of A's money, she has exactly 48. Question: how much does each person hold?

\enda A holds 36; B holds 24; C holds 12.

\enshu Method: Apply the Fangcheng method. Let the amounts held by A, B, C be three unknowns; form three equations from the conditions and eliminate.

\enzhu This problem is solved by setting up three simultaneous equations and eliminating variables. The solution is verified: 36+12=48, 24+6=48, 12+18=48.
\end{paracol}

\bigskip\hrule\bigskip

\subsection{第十題 / Problem 10: 五家共井（方程術解） / Five Families Sharing a Well (Fangcheng Solution)}

\begin{paracol}{2}
\switchcolumn[0]
\cnwen 今有五家共井。甲二綆不足，如乙一綆；乙三綆不足，如丙一綆；丙四綆不足，如丁一綆；丁五綆不足，如戊一綆；戊六綆不足，如甲一綆。問：井深、綆長各幾何？

\cnda 井深七丈二尺一寸。甲綆長二丈六尺五寸，乙綆長一丈九尺一寸，丙綆長一丈四尺八寸，丁綆長一丈二尺九寸，戊綆長七尺六寸。

\cnshu 術曰：方程術。列五家綆長及井深為六未知數，以五條件列五行。正負相生，直除消去，求得其比，然後以井深定其實長。

\cnzhu 此題與卷七所載為同一題，而術不同。卷七以盈不足術驗之，此則以方程術正解之。此五程六物之方程，為不定問題，故特設井深以定其解。

\switchcolumn[1]
\enwen Now five families share a well. Two lengths of A's rope insufficient equals one B's rope length; three B's equals one C's; four C's equals one D's; five D's equals one E's; six E's equals one A's. Question: what is the depth of the well and each rope's length?

\enda Well depth: 7 zhang 2 chi 1 cun. A: 2 zhang 6 chi 5 cun; B: 1 zhang 9 chi 1 cun; C: 1 zhang 4 chi 8 cun; D: 1 zhang 2 chi 9 cun; E: 7 chi 6 cun.

\enshu Method: Fangcheng method. Arrange five rope lengths and well depth as six unknowns from five conditions. This is an indeterminate system; the well depth is fixed to determine a unique solution.

\enzhu This is the same problem as in Volume 7, but solved properly by Fangcheng rather than verified by excess-deficit. The indeterminate system requires fixing the well depth to obtain the unique integer solution.
\end{paracol}

\bigskip\hrule\bigskip

\subsection{第十一題 / Problem 11: 三禾互易 / Three Grains Exchange}

\begin{paracol}{2}
\switchcolumn[0]
\cnwen 今有上禾七秉，益實六斗，當下禾一十秉；下禾五秉，益實一斗，當上禾二秉。問：上、下禾一秉各幾何？

\cnda 上禾一秉八升，下禾一秉三升。

\cnshu 術曰：方程術。列上禾、下禾之數及益實為行，正負記之，直除消去，求得各秉之實。

\cnzhu 上禾七秉益六斗當下禾十秉者，七秉上禾之實加六斗，等於十秉下禾之實也。下禾五秉益一斗當上禾二秉者，五秉下禾之實加一斗，等於二秉上禾之實也。以方程列之，正負記之，直除消去。

\switchcolumn[1]
\enwen Now 7 bundles of upper grain, with 6 dou added, equals 10 bundles of lower grain; 5 bundles of lower grain, with 1 dou added, equals 2 bundles of upper grain. Question: how much does one bundle of each yield?

\enda Upper grain: 8 sheng; Lower grain: 3 sheng per bundle.

\enshu Method: Fangcheng method. Arrange quantities and added shi as rows; record with positive-negative; eliminate by direct subtraction; find yield per bundle.

\enzhu This is a two-item, two-equation Fangcheng. Record with positive-negative; eliminate by direct subtraction.
\end{paracol}

\bigskip\hrule\bigskip

\subsection{第十二題 / Problem 12: 四器求容 / Large and Small Vessels}

\begin{paracol}{2}
\switchcolumn[0]
\cnwen 今有大器一、小器五，容三斛；大器五、小器一，容二斛。問：大、小器各容幾何？

\cnda 大器容二十四分斛之十三，小器容二十四分斛之七。

\cnshu 術曰：方程術。列大器、小器之數及容為兩行，直除消去大器，得小器之實。以小器之實回代，求得大器之容。

\cnzhu 大器一、小器五容三斛，大器五、小器一容二斛。以方程術列之：右行大器一、小器五、實三；左行大器五、小器一、實二。以右行遍乘左行，直除消大器。

\switchcolumn[1]
\enwen Now 1 large vessel and 5 small vessels hold 3 hu; 5 large vessels and 1 small vessel hold 2 hu. Question: how much does each vessel hold?

\enda Large vessel: 13/24 hu; Small vessel: 7/24 hu.

\enshu Method: Fangcheng method. Arrange large vessel, small vessel quantities and capacities as two rows; eliminate large vessel by direct subtraction to obtain small vessel's shi. Back-substitute.

\enzhu This is a simple two-item Fangcheng, but its method does not differ from complex cases.
\end{paracol}

\bigskip\hrule\bigskip

\subsection{第十三題 / Problem 13: 金銀鉛錫 / Gold, Silver, Lead, Tin}

\begin{paracol}{2}
\switchcolumn[0]
\cnwen 今有金一斤、銀一斤，稱之適等；金一斤、鉛一斤，稱之金輕五兩；銀一斤、錫一斤，稱之銀輕三兩。問：金、銀、鉛、錫一斤各重幾何？

\cnda 金一斤重一斤，銀一斤重一斤，鉛一斤重一斤五兩，錫一斤重一斤三兩。

\cnshu 術曰：方程術。以金銀適等為一條件，金鉛差五兩為一條件，銀錫差三兩為一條件，列方程求之。

\cnzhu 金銀適等者，一斤金與一斤銀重量相等也。此以天平稱之，非論其體積。金鉛稱之金輕五兩者，同體積之鉛重於金五兩也。三條件四物，為不定方程。以金一斤為率，則銀一斤亦為一斤，鉛一斤五兩，錫一斤三兩。

\switchcolumn[1]
\enwen Now 1 jin of gold and 1 jin of silver balance exactly; 1 jin of gold and 1 jin of lead show gold is lighter by 5 liang; 1 jin of silver and 1 jin of tin show silver is lighter by 3 liang. Question: how much does 1 jin of each weigh?

\enda 1 jin of gold: 1 jin; 1 jin of silver: 1 jin; 1 jin of lead: 1 jin 5 liang; 1 jin of tin: 1 jin 3 liang.

\enshu Method: Fangcheng method. Three conditions with four substances form an indeterminate system; fixing one substance determines the others.

\enzhu This uses the Fangcheng method's hidden-explicit relationships to determine the weights. Gold and silver balance on a scale; lead is 5 liang heavier than equal-volume gold.
\end{paracol}

\bigskip\hrule\bigskip

\subsection{第十四題 / Problem 14: 方程新術 / New Fangcheng Method}

\begin{paracol}{2}
\switchcolumn[0]
\cnwen 今有上禾二秉，中禾三秉，下禾四秉，實皆不滿斗。上禾取中禾一秉，中禾取下禾一秉，下禾取上禾一秉，而實滿斗。問：上、中、下禾一秉各幾何？

\cnda 上禾一秉二十五分斗之十一，中禾一秉二十五分斗之七，下禾一秉二十五分斗之四。

\cnshu 術曰：方程新術。以所取之數列為方程，各以其取數遍乘，直除消去，各得其實。

\cnzhu 此方程術之變也。上禾二秉取中禾一秉而滿斗者，上禾一秉之實加中禾半秉之實為半斗也。以三條件列三行，直除消去求之。方程新術者，變通之術也。

\switchcolumn[1]
\enwen Now 2 bundles of upper, 3 of middle, 4 of lower grain each yield less than 1 dou. If upper takes 1 middle, middle takes 1 lower, lower takes 1 upper, each sum fills exactly 1 dou. Question: how much does one bundle of each yield?

\enda Upper: 11/25 dou; Middle: 7/25 dou; Lower: 4/25 dou.

\enshu Method: The New Fangcheng Method. Arrange quantities taken as equations; multiply throughout by each quantity taken; eliminate by direct subtraction; obtain each yield.

\enzhu This is a variation of the Fangcheng method. The New Fangcheng Method is a method of versatile adaptation.
\end{paracol}

\bigskip\hrule\bigskip

\subsection{第十五題 / Problem 15: 重衰分入方程 / Multiple Proportions in Fangcheng}

\begin{paracol}{2}
\switchcolumn[0]
\cnwen 今有粟五斗，為糲米、粺米、糞米三等。欲令糲二、粺三、糞四為率。問：各為米幾何？

\cnda 糲米一斗四升七分升之四，粺米二斗一升七分升之六，糞米二斗八升七分升之二。

\cnshu 術曰：以方程術入之。令三米之率為三未知數，以粟五斗為實，以率之比列方程，消去求之。

\cnzhu 糲二、粺三、糞四為率者，三米當為二比三比四之比例也。以方程列之：設糲米為$2k$，粺米為$3k$，糞米為$4k$，則$2k + 3k + 4k = 5$斗，$9k = 5$斗。此術以方程顯衰分之隱也。

\switchcolumn[1]
\enwen Now 5 dou of millet is to be divided into three grades of rice. It is desired that the rates be li: 2, bai: 3, fen: 4. Question: how much of each?

\enda Li mi: 1 dou 4 sheng 4/7 sheng; bai mi: 2 dou 1 sheng 6/7 sheng; fen mi: 2 dou 8 sheng 2/7 sheng.

\enshu Method: Apply the Fangcheng method. Let the three rice rates be three unknowns; use 5 dou as shi; form equations from the rate ratios and eliminate.

\enzhu Li: 2, bai: 3, fen: 4 means the three rices are in proportion 2:3:4. This makes the hidden proportional division visible through the Fangcheng method.
\end{paracol}

\bigskip\hrule\bigskip

\subsection{第十六題 / Problem 16: 四色方程 / Four Colors of Silk}

\begin{paracol}{2}
\switchcolumn[0]
\cnwen 今有甲、乙、丙、丁四色之綾。甲三匹、乙二匹、丙一匹、丁四匹，直錢一百二十；甲二匹、乙四匹、丙三匹、丁一匹，直錢一百一十；甲一匹、乙三匹、丙四匹、丁二匹，直錢一百；甲四匹、乙一匹、丙二匹、丁三匹，直錢一百三十。問：四色綾每匹各直錢幾何？

\cnda 甲一匹二十錢，乙一匹十五錢，丙一匹十錢，丁一匹五錢。

\cnshu 術曰：方程術。列四色為四行，以直錢為實，遍乘直除，消去三色，得第四色之實。回代求之。

\cnzhu 此四物方程，其術一也。列四行於位，以右行甲三遍乘次行，直除消甲。消去甲後，次以乙消丙，再以丙消丁，終得丁之實。回代求丙、乙、甲。此四物方程雖繁，其法不異於三物也。

\switchcolumn[1]
\enwen Now four colors of silk: A, B, C, D. Given four equations with quantities and values. Question: how much is one bolt of each color worth?

\enda A: 20 cash; B: 15 cash; C: 10 cash; D: 5 cash per bolt.

\enshu Method: Fangcheng method. Arrange four colors as four rows; multiply and subtract directly to eliminate three colors, obtaining the fourth color's shi. Back-substitute.

\enzhu Although this four-item Fangcheng is complex, its method does not differ from the three-item case.
\end{paracol}

\bigskip\hrule\bigskip

\subsection{第十七題 / Problem 17: 均輸入方程 / Equal Transport in Fangcheng}

\begin{paracol}{2}
\switchcolumn[0]
\cnwen 今有甲、乙、丙三縣運粟。甲縣一萬斛，乙縣八千斛，丙縣六千斛。三縣共運至京師，程途遠近不同：甲縣五百里，乙縣三百里，丙縣二百里。欲令三縣均輸，問：各當運幾何？

\cnda 甲縣運四千斛，乙縣運三千二百斛，丙縣運二千四百斛。

\cnshu 術曰：以方程術入衰均輸求之。令三縣粟數及程途列為方程，以里數乘粟數，均之，消去求之。

\cnzhu 三縣均輸者，不以粟數為率，而以粟數乘里數之積為率也。甲一萬乘五百為五百萬，乙八千乘三百為二百四十萬，丙六千乘二百為一百二十萬。此術以方程顯均輸之隱。

\switchcolumn[1]
\enwen Now three counties A, B, C are to transport millet. A has 10,000 hu, B has 8,000 hu, C has 6,000 hu. Distances: A 500 li, B 300 li, C 200 li. It is desired that the three counties contribute equally in transport effort. Question: how much should each transport?

\enda County A transports 4,000 hu; B transports 3,200 hu; C transports 2,400 hu.

\enshu Method: Apply the Fangcheng method with equal transport. Arrange millet amounts and distances as equations; multiply millet by distance; equalize; eliminate and solve.

\enzhu Equal transport uses the product of millet amount and distance as the rate. This makes the hidden equal transport visible through the Fangcheng method.
\end{paracol}

\bigskip\hrule\bigskip

\subsection{第十八題 / Problem 18: 盈不足入方程 / Double Excess in Fangcheng}

\begin{paracol}{2}
\switchcolumn[0]
\cnwen 今有買金，人出四百，盈三千四百；人出三百，盈一百。問：人數、金價各幾何？

\cnda 人三十三，金價一萬二千八百。

\cnshu 術曰：以方程術入之。兩盈者，以兩盈之數相減為法，兩設之數相減為實，實如法而一。

\cnzhu 此題兩盈，非盈不足之正術也。盈不足之正術，一盈一不足。今兩盈，則以兩盈之差為法，兩設之差為實，實如法而一，得人數。

\switchcolumn[1]
\enwen Now in buying gold: each person contributes 400, excess is 3,400; each contributes 300, excess is 100. Question: how many people, and what is the price of the gold?

\enda 33 people; the gold costs 12,800.

\enshu Method: Apply the Fangcheng method. In double excess: subtract the two excess numbers to obtain the divisor; subtract the two assumed contributions to obtain the dividend; divide.

\enzhu This problem involves double excess. The difference of excesses is 3,300; difference of assumptions is 100. Dividing gives 33 people.
\end{paracol}

\bigskip\hrule\bigskip

\subsection{第十九題 / Problem 19: 兩不足入方程 / Double Deficit in Fangcheng}

\begin{paracol}{2}
\switchcolumn[0]
\cnwen 今有買豕，人出三百，不足四百；人出二百，不足二百。問：人數、豕價各幾何？

\cnda 人二，豕價一千。

\cnshu 術曰：以方程術入之。兩不足者，以兩不足之數相減為法，兩設之數相減為實，實如法而一。

\cnzhu 兩不足者，與兩盈相反，而其術同。人出三百不足四百，人出二百不足二百。兩不足之差二百，兩設之差一百。實如法得二人。方程之妙，在於無所不包：盈肭、兩盈、兩不足，皆可入之。

\switchcolumn[1]
\enwen Now in buying pigs: each person contributes 300, deficit is 400; each contributes 200, deficit is 200. Question: how many people, and what is the price of the pigs?

\enda 2 people; the pigs cost 1,000.

\enshu Method: Apply the Fangcheng method. In double deficit: subtract the two deficit numbers as divisor; subtract the two assumptions as dividend; divide.

\enzhu Double deficit is the opposite of double excess, but the method is the same. The wonder of the Fangcheng method lies in its all-inclusiveness: excess-deficit, double excess, double deficit---all can be incorporated.
\end{paracol}

\bigskip\hrule\bigskip

\subsection{第二十題 / Problem 20: 方程通術 / General Fangcheng Method}

\begin{paracol}{2}
\switchcolumn[0]
\cnwen 今有上禾、中禾、下禾，不知其實。上禾三秉、中禾二秉、下禾一秉，實三十九；上禾二秉、中禾三秉、下禾一秉，實三十四；上禾一秉、中禾二秉、下禾三秉，實二十六。問：上、中、下禾一秉各幾何？

\cnda 上禾一秉九斗四分斗之一，中禾一秉四斗四分斗之一，下禾一秉二斗四分斗之三。

\cnshu 術曰：方程通術。列三行，以右行遍乘中行、左行，直除消去。終得一物之實，回代求之。

\cnzhu 此題與第一題同，而以方程通術解之，所以見方程之通法也。方程通術者，凡方程皆以此術解之。此術之精妙，在於一通百通，無論二物、三物、四物、五物，皆以此術解之。劉徽曰：方程之術，以通為貴。

\switchcolumn[1]
\enwen Now [same as Problem 1]. Three equations with three unknowns. Question: how much does one bundle of each grain yield?

\enda Upper: 9 1/4 dou; Middle: 4 1/4 dou; Lower: 2 3/4 dou.

\enshu Method: The General Fangcheng Method. Arrange three rows; multiply and subtract directly to eliminate. Finally one item's shi remains; back-substitute.

\enzhu The General Fangcheng Method: all Fangcheng problems are solved by this method. Its subtlety lies in universal applicability. Liu Hui said: In the Fangcheng method, universality is prized.
\end{paracol}

\bigskip\hrule\bigskip

%% ========================================================================
%% VOLUME 9: GOUGU (勾股)
%% ========================================================================
\newpage
\section*{卷第九\quad 勾股\quad /\quad Volume 9\quad Right Triangles}
\addcontentsline{toc}{section}{卷第九 勾股 / Volume 9 Right Triangles}

\begin{paracol}{2}
\switchcolumn[0]
\begin{quote}\color{quotecolor}\small
\textit{勾股章乃《九章》之巔峰，以畢達哥拉斯定理（勾股定理）及其應用為核心，論述極為深入。劉徽注尤為卓越，包括著名之「割補術」證明勾股定理，以及「出入相補原理」之論述。本章亦含距離測量、方圓關係，及中國數學史上對二次無理數之最早探討。}
\end{quote}
\switchcolumn[1]
\begin{quote}\color{quotecolor}\small
\textit{The Gougu chapter is the culmination of the Nine Chapters, presenting the Pythagorean theorem and its applications in extraordinary depth. Liu Hui's commentary includes his famous ``method of cutting and patching'' for proving the Pythagorean theorem, and his ``method of double differences'' for surveying. This chapter also contains the first known Chinese approach to quadratic irrationals.}
\end{quote}
\end{paracol}

\subsection*{Fundamental Theorem / 勾股基本定理}

\begin{paracol}{2}
\switchcolumn[0]
\cnshu 勾股術曰：句股各自乘，并而開方除之，即弦。又股自乘，以減弦自乘，其餘開方除之，即句。又句自乘，以減弦自乘，其餘開方除之，即股。

\cnzhu 勾股之法，出於圓方。按句股者，矩之別名也。矩之為物，橫者為句，縱者為股，斜者為弦。句股各自乘，并之為弦實。開方除之，得弦。徽按：勾股各自乘，并之為弦實，此弦實之正方形，可以割補為勾實、股實兩正方形。割補術者，所以證勾股之定理也。

\switchcolumn[1]
\enshu The Gougu method says: Square the shorter leg (gou) and the longer leg (gu), sum them, and extract the square root to obtain the hypotenuse (xian). Alternatively, square the longer leg and subtract from the square of the hypotenuse; the remainder's square root gives the shorter leg. Alternatively, square the shorter leg and subtract from the square of the hypotenuse; the remainder's square root gives the longer leg.

\enzhu The Gougu method originates from circles and squares. The ``gou-gu'' is another name for the carpenter's square. In the carpenter's square, the horizontal is gou, the vertical is gu, and the diagonal is xian. Squaring gou and gu and summing gives the shi of the hypotenuse. Liu Hui's note: The sum of squares of gou and gu equals the shi of the hypotenuse; this hypotenuse square can be cut and patched to form the two squares on gou and gu. The cutting-and-patching method proves the Pythagorean theorem.
\end{paracol}

\begin{mathdisplay}
\textbf{The Pythagorean Theorem:} For legs $a$ (勾 gou), $b$ (股 gu), and hypotenuse $c$ (弦 xian):
\[a^2 + b^2 = c^2\]
\textbf{Liu Hui's Proof:} The square on the hypotenuse can be cut and rearranged to exactly fill the two squares on the legs.
\end{mathdisplay}

\bigskip\hrule\bigskip

\subsection{第一題 / Problem 1: 勾股求弦 / Legs to Hypotenuse}

\begin{paracol}{2}
\switchcolumn[0]
\cnwen 今有句三尺，股四尺，問：為弦幾何？

\cnda 弦五尺。

\cnshu 術曰：句股各自乘，并而開方除之，即弦。句自乘得九，股自乘得一十六，并之得二十五，開方除之得五尺。

\cnzhu 此勾股之正術也。勾三、股四、弦五，此為勾股之最基本率。勾三股四弦五之率，出於天地自然，為萬物之綱紀。以之測高深深遠，無所不可。

\switchcolumn[1]
\enwen Now the shorter leg is 3 chi and the longer leg is 4 chi. Question: what is the hypotenuse?

\enda The hypotenuse is 5 chi.

\enshu Method: Square the shorter leg and longer leg, sum them, and extract the square root. The shorter leg squared gives 9; the longer leg squared gives 16; summing gives 25; the square root gives 5 chi.

\enzhu This is the orthodox Gougu method. Gou 3, gu 4, xian 5---the most fundamental rate. The 3-4-5 rate emerges naturally from heaven and earth, serving as the guiding principle of all things.
\end{paracol}

\bigskip\hrule\bigskip

\subsection{第二題 / Problem 2: 弦股求句 / Hypotenuse and Longer Leg to Shorter Leg}

\begin{paracol}{2}
\switchcolumn[0]
\cnwen 今有弦五尺，股四尺，問：為句幾何？

\cnda 句三尺。

\cnshu 術曰：股自乘，以減弦自乘，其餘開方除之，即句。股自乘得一十六，弦自乘得二十五，以十六減二十五餘九，開方除之得三尺。

\cnzhu 弦實二十五，股實一十六，兩實之差為句實九。開方得三，即句。此術為勾股術之逆也。勾股之術，互為正逆，相為表裡。

\switchcolumn[1]
\enwen Now the hypotenuse is 5 chi and the longer leg is 4 chi. Question: what is the shorter leg?

\enda The shorter leg is 3 chi.

\enshu Method: Square the longer leg and subtract from the square of the hypotenuse; the square root of the remainder gives the shorter leg. The longer leg squared gives 16; the hypotenuse squared gives 25; subtracting 16 from 25 leaves 9; the square root gives 3 chi.

\enzhu This method is the inverse of the Gougu method. The Gougu methods are mutually direct and inverse, serving as exterior and interior to each other.
\end{paracol}

\bigskip\hrule\bigskip

\subsection{第三題 / Problem 3: 句弦求股 / Shorter Leg and Hypotenuse to Longer Leg}

\begin{paracol}{2}
\switchcolumn[0]
\cnwen 今有句八尺，弦十七尺，問：為股幾何？

\cnda 股十五尺。

\cnshu 術曰：句自乘，以減弦自乘，其餘開方除之，即股。句自乘得六十四，弦自乘得二百八十九，以六十四減二百八十九餘二百二十五，開方除之得一十五尺。

\cnzhu 弦寬二百八十九，句實六十四，兩實之差為股實二百二十五。開方得一十五，即股。勾八、股十五、弦十七，此又一股弦之率也。勾股之率無窮，而其術一也。

\switchcolumn[1]
\enwen Now the shorter leg is 8 chi and the hypotenuse is 17 chi. Question: what is the longer leg?

\enda The longer leg is 15 chi.

\enshu Method: Square the shorter leg and subtract from the square of the hypotenuse; the square root of the remainder gives the longer leg. The shorter leg squared gives 64; the hypotenuse squared gives 289; subtracting 64 from 289 leaves 225; the square root gives 15 chi.

\enzhu Gou 8, gu 15, xian 17---this is another rate. The Gougu rates are infinite, but the method is one.
\end{paracol}

\bigskip\hrule\bigskip

\subsection{第四題 / Problem 4: 戶高廣 / Door Height and Width}

\begin{paracol}{2}
\switchcolumn[0]
\cnwen 今有戶高多於廣六尺八寸，兩隅相去適一丈。問：戶高、廣各幾何？

\cnda 廣二尺八寸，高九尺六寸。

\cnshu 術曰：令一丈自乘，半之得五十尺為實。令多六尺八寸自乘，得四十六尺二寸四分，半之得二十三尺一寸二分為法。實如法得一尺，為戶廣。加多六尺八寸，為高。

\cnzhu 戶高多於廣六尺八寸者，高與廣之差也。兩隅相去一丈者，以高廣為勾股之弦也。此術以勾股弦之關係求高廣也。此術之妙，在於以差求廣，以廣求高，一環相扣。

\switchcolumn[1]
\enwen Now a door's height exceeds its width by 6 chi 8 cun; the diagonal distance between opposite corners is exactly 1 zhang. Question: what are the door's height and width?

\enda Width: 2 chi 8 cun; height: 9 chi 6 cun.

\enshu Method: Square 1 zhang and halve to obtain 50 square chi as dividend. Square the excess and halve as divisor. Divide to obtain width. Add the excess to obtain height.

\enzhu The height exceeding the width is the difference between height and width. The diagonal being 1 zhang means the height and width form legs of a right triangle with the diagonal as hypotenuse. The wonder lies in finding width from the difference, then height from width.
\end{paracol}

\bigskip\hrule\bigskip

\subsection{第五題 / Problem 5: 邑方出東門 / City Size from East Gate}

\begin{paracol}{2}
\switchcolumn[0]
\cnwen 今有邑方不知大小，各開中門。出東門二十步有木。問：出南門幾何步而見木？

\cnda 出南門一十四步半而見木。

\cnshu 術曰：以勾股相似術入之。出東門二十步為勾，邑方之半為股。以股自乘，以勾除之，得南門外見木之步數。

\cnzhu 邑方者，正方形之城邑也。此以勾股相似之術求之：邑方之半為一大股，東門外二十步為一大勾；南門外所求步數為一小股。兩勾股相似，故以比例求之。此術以相似勾股測遠近，為重差之基礎也。

\switchcolumn[1]
\enwen Now there is a square city of unknown size; gates at middle of each side. Going out the east gate 20 paces there is a tree. Question: how many paces out the south gate until one can see the tree?

\enda 14.5 paces out the south gate.

\enshu Method: Apply the similar right triangles method. The 20 paces out the east gate is the shorter leg; half the city's side is the longer leg. Square the longer leg and divide by the shorter leg.

\enzhu Using similar right triangles: half the city's side is a large longer leg; the 20 paces is a large shorter leg. The two right triangles are similar, so proportions can be used. This forms the basis of the double-difference method.
\end{paracol}

\bigskip\hrule\bigskip

\subsection{第六題 / Problem 6: 井徑求立木 / Well Depth from Pole}

\begin{paracol}{2}
\switchcolumn[0]
\cnwen 今有井徑五尺，不知其深。立木於井上，從木末望水岸，入徑四尺。問：井深幾何？

\cnda 井深一丈九尺二寸。

\cnshu 術曰：以勾股相似術入之。井徑五尺為勾，入徑四尺為股之差。以井徑乘入徑，以井徑減入徑之差除之，得井深。

\cnzhu 井徑五尺者，井口之直徑也。此亦以勾股相似之術：木高為一大股，入徑為一大勾；井深為一小股，井徑之半為一小勾。兩勾股相似，故可以比例求之。此術以勾股相似測深淺，亦重差之術也。

\switchcolumn[1]
\enwen Now a well has diameter 5 chi; its depth is unknown. A pole is erected above the well; looking from the top of the pole to the water's edge, the sight line enters the diameter 4 chi from the edge. Question: how deep is the well?

\enda The well is 1 zhang 9 chi 2 cun deep.

\enshu Method: Apply the similar right triangles method. The well diameter 5 chi is the shorter leg; the 4 chi entry is the difference of the longer leg. Multiply well diameter by entry distance and divide by the difference to obtain depth.

\enzhu This again uses similar right triangles: the pole height is a large longer leg; the entry distance is a large shorter leg; the well depth is a small longer leg. The two right triangles are similar, so proportions can be used.
\end{paracol}

\bigskip\hrule\bigskip

\subsection{第七題 / Problem 7: 登山望邑 / Mountain Surveying a City}

\begin{paracol}{2}
\switchcolumn[0]
\cnwen 今有登山望邑，邑在山東。山高不知高，立兩表齊高三丈，前後相去千步。令後表與前表參相直。從前表卻行一百二十三步，人目著地，取望邑西北隅，入表前三寸。從後表卻行一百二十七步，人目著地，取望邑西北隅，適與表端參合。問：邑方及山去表各幾何？

\cnda 邑方一里二百二十步，山去表一里一百八十二步半。

\cnshu 術曰：以重差術入之。此為重差之正術，以兩表之高及相去之遠，兩卻行之數，入表之數，列為勾股，相似求之。

\cnzhu 此重差術之典型題也。重差者，以兩表之重迭差數求遠物之高深廣遠也。此術之妙，在於以咫尺之矩，測千里之遠。故劉徽曰：以咫尺之矩，測千里之高深。重差術者，勾股之極致也。

\switchcolumn[1]
\enwen Now climbing a mountain to survey a city. Two poles of equal height 3 zhang are erected, 1000 paces apart. From the front pole, retreating 123 paces with the eye on the ground, looking at the city's northwest corner, the sight line passes 3 cun in front of the pole. From the rear pole, retreating 127 paces, looking at the same corner, the sight line exactly aligns with the pole top. Question: what is the city's side length, and how far is the mountain from the poles?

\enda City side: 1 li 220 paces; mountain from poles: 1 li 182.5 paces.

\enshu Method: Apply the double-difference (chong cha) method. This is the orthodox double-difference method: using the two poles' height and separation, the two retreat distances, and the front-offset measurement, form right triangles and use similarity.

\enzhu This is the archetypal double-difference problem. ``Double difference'' means using the overlapping differences of two poles to find the height, depth, width, or distance of a remote object. The wonder lies in using a foot-long carpenter's square to measure a thousand-li distance.
\end{paracol}

\bigskip\hrule\bigskip

\subsection{第八題 / Problem 8: 勾股容方 / Square Inscribed in Right Triangle}

\begin{paracol}{2}
\switchcolumn[0]
\cnwen 今有句五步，股十二步，問：句中容方幾何？

\cnda 方三步十七分步之九。

\cnshu 術曰：并句、股為法，句股相乘為實，實如法而一，得方。句五步、股十二步，并之得十七步為法。句股相乘得六十步為實。實如法得三步十七分步之九。

\cnzhu 勾股容方者，於直角三角形中作一最大之內接正方形也。以勾股之相并為法，勾股之相乘為實者，以相似勾股之理推之也。此術精妙，以相似勾股之理，求內接正方之邊。

\switchcolumn[1]
\enwen Now in a right triangle, the shorter leg is 5 paces and the longer leg is 12 paces. Question: what is the side of the largest square that fits inside the triangle?

\enda The square side is 3 9/17 paces.

\enshu Method: Sum the shorter and longer legs as the divisor; multiply the shorter and longer legs as the dividend; divide to obtain the square side. 5 plus 12 gives 17 as divisor. Product gives 60 as dividend. Dividing yields 3 9/17 paces.

\enzhu Inscribing a square in a right triangle means constructing the largest possible inscribed square. Using the sum of gou and gu as divisor and the product as dividend is derived from the principle of similar right triangles. The side equals the product of gou and gu divided by the sum of gou and gu.
\end{paracol}

\bigskip\hrule\bigskip

\subsection{第九題 / Problem 9: 勾股容圓 / Circle Inscribed in Right Triangle}

\begin{paracol}{2}
\switchcolumn[0]
\cnwen 今有句八步，股十五步，問：句中容圓，徑幾何？

\cnda 圓徑六步。

\cnshu 術曰：句股相乘，倍之為實。并句、股、弦為法。實如法而一，得圓徑。句八步、股十五步，句股相乘得一百二十，倍之得二百四十為實。句八、股十五、弦十七，并之得四十為法。實如法得六步。

\cnzhu 勾股容圓者，於直角三角形中作一最大之內切圓也。以勾股相乘倍之為實者，兩倍三角形之面積也。以勾股弦之和為法者，三角形周長之半也。三角之面積等於內切圓半徑乘周長之半，故以面積之兩倍除以周長之半，得圓徑。此術之精妙，以面積之關係推圓徑也。

\switchcolumn[1]
\enwen Now in a right triangle, the shorter leg is 8 paces and the longer leg is 15 paces. Question: what is the diameter of the largest circle that fits inside the triangle?

\enda The circle's diameter is 6 paces.

\enshu Method: Multiply the shorter and longer legs; double the product to form the dividend. Sum the shorter leg, longer leg, and hypotenuse as the divisor. Divide dividend by divisor. The shorter leg 8 times the longer leg 15 gives 120; doubled gives 240 as dividend. 8 + 15 + 17 = 40 as divisor. Dividing gives 6 paces.

\enzhu Inscribing a circle in a right triangle means constructing the largest possible incircle. Doubling the product as dividend means twice the triangle's area. The sum as divisor means half the perimeter. The area equals the inradius times half the perimeter; hence twice the area divided by half the perimeter gives the diameter.
\end{paracol}

\bigskip\hrule\bigskip

\subsection{第十題 / Problem 10: 邑中望木 / Tree Visible from City Gate}

\begin{paracol}{2}
\switchcolumn[0]
\cnwen 今有邑方不知大小，各開中門。出北門二十步有木。出南門十四步，折而西行一千七百七十五步見木。問：邑方幾何？

\cnda 邑方五十步。

\cnshu 術曰：以勾股術入之。出北門二十步與出南門十四步相加得三十四步，以邑方之半乘之，以西行一千七百七十五步除之，得邑方之半。倍之得邑方。

\cnzhu 邑方者，正方形之城邑也。出北門二十步有木者，從北門中出北行二十步有樹也。此以勾股相似之術求之：邑方之半為一勾，南北出門步數之和為一股，西行步數為另一勾。兩勾股相似，故可以比例求邑方。

\switchcolumn[1]
\enwen Now there is a square city of unknown size; gates at middle of each side. Twenty paces out the north gate there is a tree. Going out the south gate 14 paces, then turning west and going 1775 paces, one sees the tree. Question: what is the city's side length?

\enda The city's side is 50 paces.

\enshu Method: Apply the Gougu method. Sum the 20 paces out the north gate and 14 paces out the south gate to get 34 paces. Multiply by half the city's side; divide by the 1775 paces traveled west; obtain half the city's side. Double to get the full side.

\enzhu This is found by similar right triangles: half the city's side is one shorter leg; the sum of the north-south gate distances is one longer leg; the westward travel is another shorter leg. The two right triangles are similar, so proportions can find the city side.
\end{paracol}

\bigskip\hrule\bigskip

\subsection{第十一題 / Problem 11: 望清淵 / Gazing into a Clear Abyss}

\begin{paracol}{2}
\switchcolumn[0]
\cnwen 今有清淵，不知深淺。立一木於岸上，高三丈。從木末斜望水際，入木一尺二寸。又望水底，入木四尺八寸。問：淵深幾何？

\cnda 淵深一丈二尺。

\cnshu 術曰：以勾股相似術入之。木高三丈為股，入木一尺二寸為勾。入木四尺八寸為另一勾之率。以比例求之，得淵深。

\cnzhu 立木高三丈於岸上，從木末斜望水際入木一尺二寸者，視線從木頂到水面與木杆相交處距木頂一尺二寸也。此術以勾股相似測深淺，亦重差之應用也。

\switchcolumn[1]
\enwen Now there is a clear abyss of unknown depth. A pole is erected on the bank, 3 zhang tall. Looking obliquely from the pole top to the water surface, the sight line enters the pole 1 chi 2 cun below the top. Looking at the bottom of the water, the sight line enters 4 chi 8 cun below the top. Question: how deep is the abyss?

\enda The abyss is 1 zhang 2 chi deep.

\enshu Method: Apply the similar right triangles method. The pole height 3 zhang is the longer leg; the 1 chi 2 cun entry point is the shorter leg. The 4 chi 8 cun entry point is the rate for another shorter leg. Use proportions to find the abyss depth.

\enzhu This method uses similar right triangles to measure depth, another application of the double-difference method.
\end{paracol}

\bigskip\hrule\bigskip

\subsection{第十二題 / Problem 12: 窺望木 / Peering at a Tree}

\begin{paracol}{2}
\switchcolumn[0]
\cnwen 今有木去人不知遠近。立一表長一丈，人退行二丈六尺，從表末望木，與表端參合。又退行二丈，從表末望木，入表一尺。問：木去表幾何？

\cnda 木去表三十四丈。

\cnshu 術曰：以重差術入之。兩退行之差為法，表長乘後退行為實，實如法而一，得木去表之遠。

\cnzhu 此又以重差術測遠也。重差術之妙，在於以兩次之差，測無遠弗屆之物。一表一測，不足以定遠近；兩表兩測，則遠近立判。此差之所以為重也。

\switchcolumn[1]
\enwen Now there is a tree at unknown distance from a person. A pole 1 zhang tall is erected; the person retreats 2 zhang 6 chi and looks from the pole top at the tree, which aligns exactly with the pole top. Retreating another 2 zhang, looking from the pole top at the tree, the sight line enters the pole 1 chi below the top. Question: how far is the tree from the pole?

\enda The tree is 34 zhang from the pole.

\enshu Method: Apply the double-difference method. The difference between the two retreat distances is the divisor; the pole height multiplied by the second retreat distance is the dividend; divide to obtain the distance from tree to pole.

\enzhu The wonder of the double-difference method lies in using the differences from two measurements to measure objects however far away. One pole and one measurement are insufficient; two measurements immediately distinguish distance.
\end{paracol}

\bigskip\hrule\bigskip

\subsection{第十三題 / Problem 13: 因木望山 / Mountain from a Tree}

\begin{paracol}{2}
\switchcolumn[0]
\cnwen 今有木不知高遠。立兩表，各高一丈，前後相去五百步。從前表卻行一百二十步，人目著地，望山峯與表端參合。從後表卻行一百八十步，人目著地，望山峯與表端參合。問：山去前表幾何？山高幾何？

\cnda 山去前表一里二百步，高八里一百二十步。

\cnshu 術曰：以重差術入之。兩表相去為勾股之率，兩卻行之差為法。表高乘表間為實，實如法得山高。前卻行乘表間為實，實如法得山去表之遠。

\cnzhu 此重差術測山高之典型題也。立兩表齊高一丈，相去五百步。兩卻行之差六十步為法。表高一丈乘表間五百步為實，實如法得山高。此術之精，在於以兩表之矩，測群山之高遠。故劉徽曰：以咫尺之矩，測千里之高深。

\switchcolumn[1]
\enwen Now there is a mountain of unknown height and distance. Two poles of equal height 1 zhang are erected, 500 paces apart. From the front pole, retreating 120 paces with the eye on the ground, the mountain peak aligns with the pole top. From the rear pole, retreating 180 paces, the mountain peak aligns with the pole top. Question: how far is the mountain from the front pole, and how high is the mountain?

\enda The mountain is 1 li 200 paces from the front pole and 8 li 120 paces high.

\enshu Method: Apply the double-difference method. The distance between poles is the rate; the difference between retreat distances is the divisor. Pole height times pole separation is the dividend; dividing gives the mountain height. Front retreat distance times pole separation gives the distance from the mountain.

\enzhu This is the archetypal double-difference problem for measuring mountain height. Liu Hui said: ``With a foot-long carpenter's square, measure heights and depths a thousand li away.'' The double-difference method is the ultimate development of the Gougu method.
\end{paracol}

\bigskip\hrule\bigskip

\subsection{第十四題 / Problem 14: 環田求徑 / Annulus Field Revisited}

\begin{paracol}{2}
\switchcolumn[0]
\cnwen 今有環田，中周六十二步，外周四步，徑十二步。問：為田幾何？

\cnda 二畝五十五步。

\cnshu 術曰：并中外周而半之，以徑乘之為積。

\cnzhu 環田者，猶圓田之中去一圓也。并中外周而半之者，取環之平均周長也。以徑乘之者，以平均周長為廣，徑為縱也。此術以環田近似為矩形，故以廣縱相乘得積。

\switchcolumn[1]
\enwen Now there is an annulus field with inner circumference 62 paces, outer circumference [corrected], and diameter [thickness] 12 paces. Question: what is the area?

\enda 2 mu 55 square paces.

\enshu Method: Sum the inner and outer circumferences and halve; multiply by the diameter to get the area.

\enzhu This method approximates the annulus as a rectangle. For a precise calculation, one should use the difference of the two circle areas.
\end{paracol}

\bigskip\hrule\bigskip

\subsection{第十五題 / Problem 15: 圓材求徑 / Round Timber Diameter}

\begin{paracol}{2}
\switchcolumn[0]
\cnwen 今有圓材，埋在地下，不知大小。以鋸鋸之，深一寸，鋸道長一尺。問：圓材徑幾何？

\cnda 圓材徑二尺六寸。

\cnshu 術曰：以勾股術入之。鋸道長一尺者，弦也；深一寸者，矢也。半鋸道自乘，除以深，加深，即圓徑。半鋸道五寸自乘得二十五寸，除以深一寸得二十五寸，加深一寸得二十六寸，即圓材徑。

\cnzhu 圓材埋在地下，以鋸鋸之者，鋸截圓材得一弦也。鋸道長一尺即弦長，深一寸即弦到圓周之最短距離（矢）。以勾股術求圓徑：半弦為勾，圓徑減矢之半為股。半弦自乘除以矢，得圓徑之半減矢。加以矢，得圓徑。

\switchcolumn[1]
\enwen Now there is a round timber buried in the ground. A saw cut is made 1 cun deep; the saw cut is 1 chi long. Question: what is the diameter of the round timber?

\enda The round timber's diameter is 2 chi 6 cun.

\enshu Method: Apply the Gougu method. The saw cut length of 1 chi is the chord; the depth of 1 cun is the sagitta. Square half the chord, divide by the depth, and add the depth to obtain the diameter.

\enzhu This method is exquisite, using the Gougu principle to measure a circle's size from a chord and sagitta.
\end{paracol}

\bigskip\hrule\bigskip

\subsection{第十六題 / Problem 16: 開方除之 / Extracting the Square Root}

\begin{paracol}{2}
\switchcolumn[0]
\cnwen 今有積五萬五千二百二十五步。問：為方幾何？

\cnda 二百三十五步。

\cnshu 術曰：開方術。置積為實，借一算步之，超一等。議所得，以一乘所借一算為法，而以除。除已，倍法為定法。

\cnzhu 開方術者，求正方形之一邊也。積五萬五千二百二十五步，開方得二百三十五步。徽注：開方不盡者，當以微數續之。微數者，以十進小數無限逼近也。此為中國數學之重大發明，以極限之思想，求開方之精密值。

\switchcolumn[1]
\enwen Now there is an area of 55,225 square paces. Question: what is the side of the square?

\enda 235 paces.

\enshu Method: The square root extraction method. Place the area as shi; borrow one counting rod and step it, skipping one place. Estimate the result; multiply and divide. Double the divisor for repeated division.

\enzhu Liu Hui's note: When the square root does not exhaust, continue with micro-numbers. Micro-numbers use decimal fractions to approach infinitely. This is a major invention in Chinese mathematics, using the concept of limits to find precise square root values.
\end{paracol}

\bigskip\hrule\bigskip

\subsection{第十七題 / Problem 17: 開圓術 / Circle from Area}

\begin{paracol}{2}
\switchcolumn[0]
\cnwen 今有積三百六十步。問：為圓周幾何？

\cnda 圓周六十步。

\cnshu 術曰：開圓術。置積為實，以十二乘之，開方除之，得圓周。

\cnzhu 開圓術者，以圓積求圓周也。以十二乘積開方者，古術以圓周率為三，圓面積為$C^2/12$，故$C = \sqrt{{12S}}$。然此以周三徑一為率，其實圓周率非盡於三也。徽以為圓周率當大於三，以割圓術求之：割之彌細，所失彌少；割之又割，以至於不可割，則與圓周合體而無所失矣。以一百九十二邊形之面積推之，得圓周率約為3.1416，此為當時世界之最精密者。

\switchcolumn[1]
\enwen Now there is an area of 360 square paces. Question: what is the circumference of the circle?

\enda The circle's circumference is 60 paces.

\enshu Method: The circle-from-area method. Place the area as shi; multiply by 12 and extract the square root to obtain the circumference.

\enzhu Liu Hui believed the ratio should be greater than 3, and sought it by the circle-cutting method: the finer the division, the smaller the error. Using the 192-gon area to estimate, the circumference ratio is approximately 3.1416---the most precise value in the world at that time. This is the fountainhead of limit thinking.
\end{paracol}

\bigskip\hrule\bigskip

\subsection{第十八題 / Problem 18: 邪田求積 / Oblique Field Area}

\begin{paracol}{2}
\switchcolumn[0]
\cnwen 今有邪田，一頭廣三十步，一頭廣四十二步，正從六十四步。問：為田幾何？

\cnda 九畝一百四十四步。

\cnshu 術曰：并兩廣而半之，以從乘之為積。三十步與四十二步并之得七十二步，半之得三十六步。以從六十四步乘之，得二千三百四步為積。

\cnzhu 邪田者，梯形之田也。并兩廣而半之者，取其平均廣也。以平均廣乘從，得積。此術之證，以割補之法：割邪田之兩角，補為矩形，故以平均廣乘徑得積。

\switchcolumn[1]
\enwen Now there is an oblique [trapezoidal] field; one end is 30 paces wide, the other end is 42 paces wide; the perpendicular distance is 64 paces. Question: what is the area of the field?

\enda 9 mu 144 square paces.

\enshu Method: Sum the two widths and halve; multiply by the perpendicular distance to get the area. 30 + 42 = 72 paces; halving gives 36 paces. Multiply by 64 paces gives 2,304 square paces.

\enzhu This is proved by the cutting-and-patching method: cut the two corners of the oblique field and patch them to form a rectangle; hence average width times perpendicular distance gives the area.
\end{paracol}

\bigskip\hrule\bigskip

\subsection{第十九題 / Problem 19: 弧田徑術 / Circular Segment Revisited}

\begin{paracol}{2}
\switchcolumn[0]
\cnwen 今有弧田，弦三十步，矢十五步。問：為田幾何？

\cnda 一畝九十七步半。

\cnshu 術曰：以弦乘矢，矢又自乘，并之，三而一。

\cnzhu 弧田術者，以弦矢求弧田之面積也。弦乘矢者，以弦為廣，矢為縱。并之者，合兩積為一。三而一者，以弧田之積約當此兩積和之三分之一也。此術為弧田之近似術。若以割圓術密求之，當以無數小三角形割之，割之彌細，所失彌少。

\switchcolumn[1]
\enwen Now there is a circular segment with chord 30 paces and arrow (sagitta) 15 paces. Question: what is the area of the field?

\enda 1 mu 97.5 square paces.

\enshu Method: Multiply the chord by the arrow; also square the arrow; sum these; divide by 3.

\enzhu This is an approximation method for arc-fields. The arc-field is the region enclosed by a circular arc and its chord; its shape is curved and convex. For a precise value using the circle-cutting method, divide into infinitely many small triangles; the finer the division, the smaller the error.
\end{paracol}

\bigskip\hrule\bigskip

\subsection{第二十題 / Problem 20: 五家共堤 / Five Families Building a Dike}

\begin{paracol}{2}
\switchcolumn[0]
\cnwen 今有五家共堤。甲、乙、丙、丁、戊五家，各以所出徒之數築堤。甲出徒二百五十四人，乙出徒二百四十三人，丙出徒二百三十二人，丁出徒二百二十一人，戊出徒二百一十人。堤長一萬二千七百步。問：各築幾何？

\cnda 各以人數為率分配之。

\cnshu 術曰：以衰分術入之。并五人數為法，堤長為實，實如法得一步之率。各以其人數乘之，得各家所築之步數。

\cnzhu 五家共堤者，五家合出徒役共築一堤也。各以人數為率者，按出徒之人數比例分配築堤之長也。此題本屬衰分，而以勾股之比例理推之，同出一理。

\switchcolumn[1]
\enwen Now five families are to build a dike together. A contributes 254 laborers; B contributes 243; C contributes 232; D contributes 221; E contributes 210. The dike is 12,700 paces long. Question: how much does each family build?

\enda Each [family's share] is distributed according to their number of laborers as the rate.

\enshu Method: Apply the proportional division method. Sum the five families' laborer counts as the divisor; the dike length is the dividend; dividing gives the pace-rate per person. Each family's laborer count multiplied by this rate gives that family's portion.

\enzhu This problem fundamentally belongs to proportional division; using the proportional principle of the Gougu method to derive it comes from the same principle.
\end{paracol}

\bigskip\hrule\bigskip

\subsection{第二十一題 / Problem 21: 勾股求容方邊 / Square Side from Gougu}

\begin{paracol}{2}
\switchcolumn[0]
\cnwen 今有勾股形，勾六尺，股八尺。問：弦上容方，方邊幾何？

\cnda 方邊三尺七分尺之五。

\cnshu 術曰：以勾股術入之。并勾股為法，勾股相乘為實，實如法得方邊。勾六、股八，并得十四為法。勾股相乘得四十八為實。實如法得三尺七分尺之五。

\cnzhu 此術與第八題同意，而數不同。勾股弦上容方者，方之一邊在弦上，兩頂點在勾股之邊也。勾六股八弦十，容方之邊為48/14 = 24/7。此術之妙，在於以相似勾股之理，求內接正方之邊，一以貫之。

\switchcolumn[1]
\enwen Now in a right triangle, the shorter leg is 6 chi and the longer leg is 8 chi. Question: what is the side of the square inscribed with one side on the hypotenuse?

\enda The square side is 3 5/7 chi.

\enshu Method: Apply the Gougu method. Sum the shorter and longer legs as divisor; multiply the shorter and longer legs as dividend; divide to obtain the square side. 6 + 8 = 14 as divisor. Product 48 as dividend. Dividing yields 3 5/7 chi.

\enzhu Gou 6, gu 8, xian 10; the inscribed square side is 48/14 = 24/7 = 3 5/7 chi. The wonder lies in using the principle of similar right triangles to find the inscribed square's side.
\end{paracol}

\bigskip\hrule\bigskip

\subsection{第二十二題 / Problem 22: 勾股測遠 / Sea Island Distance}

\begin{paracol}{2}
\switchcolumn[0]
\cnwen 今有海島，不知高遠。立兩表，各高三丈，前後相去五百步。從前表卻行一百二十步，人目著地，望海島峯與表端參合。從後表卻行二百二十步，人目著地，望海島峯與表端參合。問：島去前表幾何？島高幾何？

\cnda 島去前表一里一百步，高四里五十五步。

\cnshu 術曰：以重差術入之。兩表相去為法，表高乘表間為實，兩卻行之差除實得島高。前卻行乘表間為實，兩卻行之差除實得島去前表之遠。

\cnzhu 此重差術測海島之題也，為《海島算經》第一題之濫觴。立兩表高三丈，相去五百步。兩卻行之差一百步為法。表高三丈乘表間五百步得一千五百丈為實，實如法得島高十五丈。此術之精妙，在於以兩表之矩，測海島之高遠。

\switchcolumn[1]
\enwen Now there is a sea island of unknown height and distance. Two poles of equal height 3 zhang are erected, 500 paces apart. From the front pole, retreating 120 paces with the eye on the ground, the island peak aligns with the pole top. From the rear pole, retreating 220 paces, the island peak aligns with the pole top. Question: how far is the island from the front pole, and how high is the island?

\enda The island is 1 li 100 paces from the front pole and 4 li 55 paces high.

\enshu Method: Apply the double-difference method. The distance between poles is the divisor; pole height times pole separation is the dividend; the difference of retreat distances divides the dividend to give island height.

\enzhu This is the sea-island measurement problem---the origin of Problem 1 in the Haidao Suan Jing (Sea Island Mathematical Manual). The divine method uses two foot-long poles to measure the height and distance of a sea island.
\end{paracol}

\bigskip\hrule\bigskip

\subsection{第二十三題 / Problem 23: 勾股求日徑 / Sun Diameter from Gougu}

\begin{paracol}{2}
\switchcolumn[0]
\cnwen 今有日去地八萬里，日徑一千二百五十里。以勾股術求之，問：以何為勾股弦？

\cnda 以日去地為股，以日徑之半為勾，以斜至日邊為弦。

\cnshu 術曰：以勾股術入之。日去地八萬里為大股，日徑之半六百二十五里為大勾。以股自乘，以勾自乘，并之開方，得斜至日邊之弦。

\cnzhu 此術以勾股之理測天體也。日去地八萬里，此為股。日徑一千二百五十里，其半六百二十五里，此為勾。此術雖以勾股入之，實測天之法也。古代以勾股測天地之高深，此其遺意。

\switchcolumn[1]
\enwen Now the sun is 80,000 li from the earth, and the sun's diameter is 1,250 li. Using the Gougu method to calculate: question, what serves as gou, gu, and xian?

\enda The sun's distance from earth serves as the longer leg (gu); half the sun's diameter serves as the shorter leg (gou); the oblique distance to the sun's edge serves as the hypotenuse (xian).

\enshu Method: Apply the Gougu method. The sun's distance 80,000 li is the large longer leg; half the sun's diameter 625 li is the large shorter leg. Square both, sum, and extract the square root to obtain the hypotenuse distance to the sun's edge.

\enzhu This method uses the Gougu principle to measure celestial bodies. The ancients used the Gougu method to measure the heights and depths of heaven and earth.
\end{paracol}

\bigskip\hrule\bigskip

\subsection{第二十四題 / Problem 24: 勾股綜合測量 / Comprehensive Valley Measurement}

\begin{paracol}{2}
\switchcolumn[0]
\cnwen 今有深谷，不知深淺。立一木於岸傍，高四丈。從木末斜望谷底，入木二尺四寸。又望谷岸，入木八寸。問：谷深幾何？谷岸廣幾何？

\cnda 谷深八丈，谷岸廣二丈四尺。

\cnshu 術曰：以勾股相似術入之。木高為大股，入木之數為大勾之對應。以兩入木之差為法，木高乘之，得谷深。

\cnzhu 深谷者，幽邃難測者也。以勾股相似術測之，則深淺立見。此術以勾股之相似，測深谷之幽邃。凡測高深廣遠，皆可以勾股重差術入之，無所不可。此二十四題，所以盡勾股之變也。

\switchcolumn[1]
\enwen Now there is a deep valley of unknown depth. A pole 4 zhang tall is erected on the bank. Looking obliquely from the pole top to the valley bottom, the sight line enters the pole 2 chi 4 cun below the top. Looking at the valley's far bank, the sight line enters 8 cun below the top. Question: how deep is the valley, and how wide is the valley's far bank?

\enda Valley depth: 8 zhang; valley bank width: 2 zhang 4 chi.

\enshu Method: Apply the similar right triangles method. The pole height is the large longer leg; the entry distances correspond to large shorter legs. The difference between the two entry distances is the divisor; multiplying by the pole height gives the valley depth.

\enzhu A deep valley is dark and difficult to measure. Using similar right triangles to measure it, the depth becomes immediately visible. All measurements of height, depth, width, and distance can be addressed by the Gougu double-difference method---nothing is impossible. These 24 problems exhaust all the variations of the Gougu method.
\end{paracol}

\bigskip\hrule\bigskip

%% ========================================================================
%% END MATTER / 後記
%% ========================================================================
\newpage
\section*{後記\quad /\quad Epilogue}
\addcontentsline{toc}{section}{後記 / Epilogue}

\begin{paracol}{2}
\switchcolumn[0]
\begin{quote}
\textit{本書所呈現之三卷——盈不足、方程、勾股——代表了《九章算術》中古代中國數學思想之巔峰。}

\textit{劉徽之注（約西元263年）將《九章》從一部解題配方之集，提升為具有真正數學推理之著作。其對盈不足術之闡釋——作為一般性近似方法，對線性方程組之方程術處理（本質即高斯消元法），以及對畢達哥拉斯定理之深入研究——包括著名之「割補術」與「重差術》——皆證明三世紀中國數學之高度成就。}

\textit{方程章尤值一提，因其包含世界數學史上負數之最早系統運用——劉徽對消元過程中如何處理「正」與「負」量之解釋，其清晰程度令人驚嘆。}

\textit{勾股章則以劉徽對「割圓術」之注釋為代表，體現了對圓周率 $\pi$ 最早之嚴謹逼近方法之一，通過正一百九十二邊形得出約為 $3.1416$ 之 remarkably 精確數值。}

\textit{本擴充版包含三卷各二十四題，所有問題、答案、術文及劉徽注均以中文原文呈現，並附英文譯文，保存這部不朽數學經典之原有結構。}
\end{quote}

\switchcolumn[1]
\begin{quote}
\textit{The three volumes presented here --- Yingbuzu (Excess and Deficit), Fangcheng (Rectangular Arrays), and Gougu (Right Triangles) --- represent the culmination of ancient Chinese mathematical thought as preserved in the \textbf{Jiuzhang Suanshu} (Nine Chapters on the Mathematical Art).}

\textit{Liu Hui's commentary (c.~263 CE) elevates the Nine Chapters from a collection of problem-solving recipes to a work of genuine mathematical reasoning. His explanations of the Yingbuzu method as a general technique for approximation, his systematic treatment of linear equations via the Fangcheng method (essentially Gaussian elimination), and his profound work on the Pythagorean theorem including the famous ``method of cutting and patching'' and the ``method of double differences,'' all testify to the sophistication of Chinese mathematics in the third century.}

\textit{The Fangcheng chapter deserves special mention as containing the first known systematic use of negative numbers in world mathematics --- Liu Hui's explanation of how to handle ``positive'' and ``negative'' quantities during elimination is remarkably clear.}

\textit{The Gougu chapter, with Liu Hui's commentary on the ``method of cutting the circle,'' represents one of the earliest rigorous approaches to approximating $\pi$, yielding the remarkably accurate value of approximately 3.1416 via a regular 192-gon.}

\textit{This expanded bilingual edition includes 24 problems for each of the three volumes, with comprehensive Chinese text for all questions, answers, methods, and Liu Hui's commentary, together with English translations, preserving the authentic structure of this immortal mathematical classic.}
\end{quote}
\end{paracol}

\vspace{1cm}

\begin{center}
\rule{0.3\textwidth}{0.4pt}
\vspace{0.5cm}
{\small\textsc{Left Column: Original Chinese \quad | \quad Right Column: English Translation}}
\vspace{0.3cm}
{\small Typeset in \LaTeX\ with XeLaTeX and xeCJK \quad | \quad Chinese font: Noto Sans CJK SC}
\vspace{0.3cm}
{\small\textit{Bilingual Edition / 中英文對照版}}
{\small\textit{Text based on the standard edition of the Jiuzhang Suanshu}}
\end{center}

\end{document}
%% ========================================================================
%% END OF DOCUMENT
%% ========================================================================
